\documentclass[letterpaper,journal]{IEEEtran}
\usepackage{amsmath,amsfonts,amssymb}
\usepackage{algorithmic}
\usepackage{algorithm}
\usepackage{array}
\usepackage{textcomp}
\usepackage{stfloats}
\usepackage{url}
\usepackage{verbatim}
\usepackage{cite}
\usepackage{amsmath}
\usepackage{booktabs}
\usepackage{graphicx}
\usepackage[caption=false,font=normalsize,labelfont=sf,textfont=sf]{subfig}
% \usepackage{caption}
\usepackage{booktabs}
\usepackage{xcolor}
\usepackage{listings}
\usepackage{multirow}
\usepackage[most]{tcolorbox}
\usepackage[dvipsnames,svgnames,x11names]{xcolor}

\hyphenation{op-tical net-works semi-conduc-tor IEEE-Xplore}
% updated with editorial comments 8/9/2021

\begin{document}

\title{SemanticAgent: A Semantics-Aware Framework for Text-to-SQL Data Synthesis}

\author{
    \IEEEauthorblockN{
        Qiang Gao\textsuperscript{1,\textdagger},
        Zhenping Li\textsuperscript{1,\textdagger},
        Anqi Zhuo\textsuperscript{1,2},
        Yingxiao Zhao\textsuperscript{1},
        Ruifang Zhao\textsuperscript{1},
        Weibo Geng\textsuperscript{1},
        and Xiaosong Li\textsuperscript{1,*}
    }

    \IEEEauthorblockA{
        \textsuperscript{1}Center of Information Research,
        Academy of Military Science, Beijing 100142, China
    }

    \IEEEauthorblockA{
        \textsuperscript{2}School of Mathematics and Physics,
        University of Science and Technology Beijing, Beijing 100083, China
    }

    \IEEEauthorblockA{
        \textsuperscript{\textdagger}These authors contributed equally.\\
        \textsuperscript{*}Corresponding author: Xiaosong Li 
        (e-mail: lxiaosong2025@163.com).
    }
}


% The paper headers
\markboth{IEEE Transactions on Knowledge and Data Engineering}%
{Gao \MakeLowercase{\textit{et al.}}: SemanticAgent for Text-to-SQL Data Synthesis}

% \IEEEpubid{0000--0000/00\$00.00~\copyright~2026 IEEE}
% Remember, if you use this you must call \IEEEpubidadjcol in the second
% column for its text to clear the IEEEpubid mark.

\maketitle

\begin{abstract}
Existing text-to-SQL synthesis pipelines rely on syntactic checks and execution feedback: they ensure executability, but not semantic correctness.
To address this gap, we propose SemanticAgent, a semantics-aware framework for text-to-SQL data synthesis.
SemanticAgent is architected around three components: an analyzer that infers a structured knowledge base from database schemas and sampled instances, a synthesizer that generates question--SQL pairs under its guidance, and a verifier that checks each candidate against the retrieved evidence.
The verifier then iteratively diagnoses each candidate and repairs the semantic errors that execution-based filtering alone would retain.
SemanticAgent reduces the human-verified semantic error rate of the synthesized corpus from [X]\% to [Y]\% relative to the strongest baseline, and models fine-tuned on its data improve execution accuracy on semantically demanding benchmarks such as BIRD.
\end{abstract}

\begin{IEEEkeywords}
Text-to-SQL, Synthetic data generation, Data augmentation, LLM agents
\end{IEEEkeywords}

\section{Introduction}\label{sec:introduction}

\begin{figure*}[!t]
  \centering
  \includegraphics[width=.9\textwidth]{case_study.pdf}
  \caption{An example of a semantically invalid but executable SQL query. The aggregation AVG(CDSCode) over a school identifier violates domain semantics, yet passes execution-based validation.}
  \label{fig:semantic_error}
\end{figure*}

Text-to-SQL supports natural-language interaction with relational databases for non-expert users, and is increasingly used in practical scenarios such as business intelligence, healthcare analytics, and domain-specific question answering~\cite{shi2025survey}. 
However, developing effective text-to-SQL models requires large-scale, high-quality question--SQL training pairs. Constructing such pairs is costly, as each instance must be grounded in the schema and validated for semantic correctness~\cite{zhang2023sciencebenchmark}. 
This challenge is particularly pronounced in domain-specific settings, where complex schemas and scarce domain expertise make it difficult to collect sufficient training data. For instance, clinical annotation required soliciting realistic questions from hospital staff, manually linking them to clinical schemas, and recruiting domain experts to verify each resulting SQL~\cite{lee2022ehrsql}. Data synthesis offers a promising alternative, reducing annotation costs and improving coverage across diverse domains. 

Consider a question--SQL pair synthesized over a school database: the query computes \texttt{AVG(CDSCode)}, where \texttt{CDSCode} is a school identifier. The pair may pass every execution check, yet averaging an identifier is domain-meaningless. Such semantic errors survive execution-based filtering and enter the training corpus~\cite{cen2025sqlfixagent} (Fig.~\ref{fig:semantic_error}). A SQL query is semantically correct with respect to its question when it faithfully expresses the question's intended operation and respects the domain knowledge of the underlying database~\cite{liu2025nl2sqlbugs}; in synthesis, the question itself must also ask something the domain supports. Existing approaches provide increasingly rich signals for text-to-SQL generation and data construction. Schema-aware methods constrain SQL generation through grammar, schema linking, and structured schema representations~\cite{yu2021grappa,li2023resdsql,liu2025}. Text-to-SQL data synthesis pipelines then use LLMs to generate candidate queries and retain those that run successfully through execution-based filtering~\cite{ye2023symgen,yang2024synthesizing,li2025omnisql,zhang2025exesql}. While these mechanisms have substantially improved syntactic validity, schema alignment, and executability, a fundamental gap remains: they excel at structural and execution checks but lack the explicit evidence required to judge whether a query is semantically correct with respect to its question~\cite{sun2023sqlpalm,cen2025sqlfixagent}.

\begin{figure*}[ht]
  \centering
  \includegraphics[width=.7\textwidth]{knowledge.pdf}
    \caption{Comparison of database context used in LLM-based text-to-SQL. DDL (B) provides structural schema information, and M-schema (C) further enriches it with schema organization and example values, but both remain schema-level representations. In contrast, our knowledge base (A) makes domain-specific semantics and constraints explicit, enabling semantic checking during both generation and refinement.}
  \label{fig:knowledge}
\end{figure*}

To address this gap, we introduce SemanticAgent, a semantics-aware synthesis framework that replaces the model's internal knowledge with explicit, retrieved evidence. Our framework is architected around three components: an analyzer, a synthesizer, and a verifier. The entire synthesis process is orchestrated through a structured, three-phase protocol: (1) Semantic Analysis, where the analyzer infers a structured knowledge base of business rules and semantic constraints from the schema and sampled instances (Fig.~\ref{fig:knowledge} and App.~\ref{app:er-analysis}); (2) Controlled Authoring, where the synthesizer composes the question, rationale, and SQL stepwise under its guidance; and (3) Diagnosis, where the verifier examines each candidate against retrieved evidence and repairs detected semantic errors.

We evaluate SemanticAgent on both fronts: the semantic correctness of the synthesized pairs and the downstream performance of models fine-tuned on them. On BIRD under Qwen2.5-Coder-7B, models fine-tuned on SemanticAgent data improve execution accuracy (EX) by 2.6 points over those trained on OmniSQL data~\cite{li2025omnisql}. On the data side, a direct human evaluation of the synthesized corpora yields a semantic error rate of [Y]\% for SemanticAgent, versus [X]\% for the strongest baseline, while diversity remains competitive.

The key contributions of this work are threefold:
\begin{itemize}
\item An Evidence-Based Framework for Text-to-SQL Data Synthesis: We propose SemanticAgent, which grounds text-to-SQL data synthesis in explicit, retrieved evidence from the database rather than the model's internal knowledge.
\item Direct Verification and Repair of Semantic Errors: SemanticAgent checks every candidate against the inferred knowledge base in an iterative diagnose--retrieve--correct loop, and the semantic errors of the synthesized corpus are measured by direct human evaluation rather than inferred from execution.
\item Improved Data Quality and Downstream Performance: Human evaluation reduces the semantic error rate of the synthesized corpus from [X]\% to [Y]\% relative to the strongest baseline, and models fine-tuned on SemanticAgent data improve execution accuracy by up to 2.6 points on BIRD over OmniSQL~\cite{li2025omnisql}.
\end{itemize}


\section{Related Work}\label{sec:related}

\subsection{Synthetic text-to-SQL Data Generation}

Manually annotating question--SQL pairs is costly and difficult to scale, especially in domain-specific settings~\cite{zhang2023sciencebenchmark,lee2022ehrsql}. Synthetic data generation has therefore emerged as a scalable alternative, and recent work demonstrates clear downstream gains across diverse domains~\cite{yang2024synthesizing,li2025omnisql}. Existing approaches improve the quality of synthesized data along two dimensions: candidate generation and quality refinement.

Early approaches synthesize candidates without LLMs: they instantiate predefined query templates across schemas~\cite{yu2021grappa}, decompose queries into abstract syntax trees and recombine reusable components~\cite{wu2021,baumgartner2024synql,guo2025sqlforge}, or derive questions from SQL through lexical substitution~\cite{yang2021hierarchical} and seq-to-seq generation~\cite{guo2018question}. More recent approaches delegate generation itself to LLMs, which produce richer and more varied question--SQL pairs through schema-aware prompting and in-context learning~\cite{ye2023symgen,yang2024synthesizing,liu2024reformer,caferoglu2025singsql}, with semantic correctness left to the model itself.

The second dimension, quality refinement, then improves the generated candidates. Ye et al.~\cite{ye2023symgen} filter candidates by execution; Yang et al.~\cite{yang2024synthesizing} add critic models and preference supervision to correct errors from weaker generators; Liu et al.~\cite{liu2024reformer} alternate generation and validation in an iterative loop; and Caferoglu et al.~\cite{caferoglu2025singsql} compare each question against a natural-language explanation of its SQL. All of them, however, rely on the same two signals: execution results and model judgments.

None of these signals examines what a query means. A query can pass every check while aggregating an identifier or applying a filter that contradicts the domain knowledge of the underlying database, the kind of error that survives into training data~\cite{liu2025nl2sqlbugs}.

\subsection{Knowledge and Reasoning Techniques}

Modern LLMs have also motivated two lines of work that improve text-to-SQL generation itself. The first enriches the context available to the model: external knowledge injected into prompts helps capture rules that the schema alone does not state~\cite{hong2024knowledge,han2024bridging}, and schema linking and retrieval methods align natural-language mentions with the right schema elements~\cite{wang2025linkalign,glass2025extractive,fan2024confidence,guo2025multipattern,zhang2025structure}. The second decomposes generation into intermediate steps: chain-of-thought prompting~\cite{wei2022chain,chang2023prompt}, reasoning decomposition~\cite{sui2025chain,dai2025parsql}, and clause-level parsing~\cite{xu2024cop,pourreza2025chasesql} expose the decisions behind a query, enabling finer-grained analysis and more targeted correction.

A common feature of both lines is that the added knowledge and reasoning serve the generation process itself: the model produces better queries, but the semantic correctness of the synthesized data is not examined.

\subsection{Error Detection and Correction}
\label{subsec:error-correction}

A parallel line of work detects and corrects errors in generated SQL queries, and the frameworks differ mainly in what they check against. The simplest reruns generation whenever the database rejects a query~\cite{wang2025macsql}. A family of self-critique methods lets the model judge itself: MAGIC~\cite{askari2025magic} distills past failures into a fixed guideline, SQLFixAgent~\cite{cen2025sqlfixagent} generates several candidate repairs and selects the most fitting one, and GBV-SQL~\cite{chen2026gbvsql} compares a natural-language translation of the query with the question. EMLC~\cite{lei2026emlc} adds fixed structural checks on the schema, the query skeleton, and executability, rewriting the query whenever a check fails. SQLENS~\cite{gong2025sqlens} pools the broadest set of checks, weighing rules derived from the database and from the model, and PV-SQL~\cite{tian2026pvsql} consults the data itself, sending probe queries to resolve ambiguities and checking the query against rules extracted from the question text.

What these frameworks share, however, is the basis of their checks: execution feedback, the model's own judgment, and fixed structural rules. While effective at catching queries that fail to execute or look structurally suspicious, none of them is grounded in per-database knowledge of what the columns and relations actually mean. Usages that execute successfully yet misuse an element therefore pass through. Error taxonomies themselves have been organized by the structure of the SQL query~\cite{liu2025nl2sqlbugs}; the taxonomy we build on in Section~\ref{subsec:refinement} instead organizes semantic errors by the evidence required to establish each class, so that each diagnosis names the knowledge needed to correct it.

This is where SemanticAgent differs. Instead of relying on execution feedback and model judgment, it grounds both diagnosis and correction in a structured semantic knowledge base built per database, and applies them where the errors matter most: at data-synthesis time, when a wrong pair becomes a training signal rather than a single failed query. The quality of this diagnosis and repair is itself measured against human annotation (Sections~\ref{subsec:diagnosis-evaluation} and~\ref{subsec:paired-refinement}).

\section{Methods}
\label{sec:methods}

We aim to synthesize text-to-SQL training data in which each SQL query is semantically correct with respect to its question, rather than merely executable. To this end, we develop \textbf{SemanticAgent}, a semantics-aware synthesis framework built on a structured semantic knowledge base $\mathcal{K}$. The framework operates in three phases. First, it infers multi-level semantic knowledge from the schema and sampled instances. Second, it generates question--SQL drafts under the guidance of $\mathcal{K}$. Finally, it diagnoses and refines semantic errors by verifying each candidate against evidence retrieved from $\mathcal{K}$. Figure~\ref{fig:overview} illustrates the overall architecture.

\begin{figure*}[!t]
 \centering
 \includegraphics[width=0.9\textwidth]{overview.pdf}
  \caption{Overview of SemanticAgent. Starting from database schemas, sampled table values, and synthesis controls, SemanticAgent (1) infers a multi-level semantic knowledge base, (2) generates question--SQL drafts under the guidance of the knowledge base, and (3) diagnoses and refines semantic errors by retrieving supporting evidence from it. The output is a semantically validated triple admitted to the synthetic corpus.}
 \label{fig:overview}
\end{figure*}

\subsection{Problem Formulation}
\label{subsec:problem-def}

Given a relational database $\mathcal{D} = (\mathcal{S}, \mathcal{I})$ with schema $\mathcal{S}$ and instances $\mathcal{I}$, we synthesize a training dataset $\mathcal{D}_{\text{syn}} = \{(q_i, s_i, r_i)\}_{i=1}^{N}$, where each sample consists of a natural language question $q_i$, an SQL query $s_i$, and a rationale $r_i$. A triple $(q, s, r)$ is semantically correct if the SQL query $s$ is syntactically valid and executable and three conditions hold: \emph{question validity}, the question $q$ asks something the domain supports; \emph{intent match}, $s$ faithfully expresses the intended operation of $q$; and \emph{usage consistency}, $s$ respects the domain knowledge of the database in how it uses each column, value, and join~\cite{liu2025nl2sqlbugs}. The first condition is specific to data synthesis: because the question is itself synthesized, it may ask something the domain does not support. The rationale $r$ must faithfully record how $s$ is derived from $q$. Because the question is itself synthesized, correctness further requires that $q$ ask something the domain supports. Since schemas alone do not encode the domain knowledge needed to judge these conditions, we infer a structured semantic knowledge base $\mathcal{K}$ from $\mathcal{S}$ and $\mathcal{I}$ (Section~\ref{subsec:semantic-construction}). Multiple semantic checks grounded in $\mathcal{K}$ hold every triple to this definition, with $\mathcal{K}$ guiding generation (Section~\ref{subsec:generation}) and supplying the evidence for diagnosis and correction (Section~\ref{subsec:refinement}). See Appendix~\ref{app:training_format} for the detailed data format.


\subsection{System Architecture}
\label{subsec:architecture}

SemanticAgent consists of three components (Figure~\ref{fig:overview}): the analyzer, the synthesizer, and the verifier. The analyzer is tasked with inferring a structured semantic knowledge base $\mathcal{K}$ from the schema and sampled instances, capturing type constraints, column semantics, and cross-table relations (Section~\ref{subsec:semantic-construction}). Guided by $\mathcal{K}$, the synthesizer composes the triples $(q, s, r)$, where the rationale $r$ records the intermediate reasoning from the question to the SQL query (Section~\ref{subsec:generation}). The verifier then examines every triple through execution, schema, and semantic checks, and repairs the detected errors (Section~\ref{subsec:refinement}).

A critical design choice in SemanticAgent is that the knowledge base $\mathcal{K}$ serves both the synthesizer and the verifier, rather than supporting generation alone. The verifier reuses both artifacts produced upstream: it retrieves corrective evidence from $\mathcal{K}$, and inspects the rationale $r$ composed by the synthesizer, to detect semantic errors and apply targeted correction while preserving valid reasoning steps. By grounding generation and verification in the same knowledge base, we ensure that the knowledge guiding composition is also the knowledge against which the finished triple is examined.

The three components execute in a fixed sequence rather than via autonomous tool selection: each component consumes the outputs of its predecessors, and the verifier internally iterates a \textsc{Diagnose}--\textsc{Retrieve}--\textsc{Correct} loop (Eq.~\ref{eq:refinement}) for at most three iterations. We justify this fixed orchestration by recognizing that the verifier's examination depends on artifacts produced upstream: without $\mathcal{K}$ there is no evidence to retrieve, and without $r$ there is no reasoning to inspect.

\subsection{Inferring Multi-Level Semantic Knowledge}
\label{subsec:semantic-construction}

Judging semantic correctness requires knowledge that the schema alone does not state, such as which columns admit aggregation and which transactions count as revenue. Much of this knowledge is recoverable from the database itself. We therefore infer a structured semantic knowledge base $\mathcal{K}$ through a six-stage hierarchical process, summarized in Table~\ref{tab:extraction-pipeline}; the process runs automatically on each target database and requires no manual rule customization.

The six stages progress from structural to semantic information: schema metadata and domain-level rules are inferred first, followed by intra-table semantics, including field types, column semantics, and table constraints, and finally cross-table relations and dependencies.

Formally, let $\mathcal{K}_t$ denote the knowledge extracted at stage $t$. At each stage, an LLM-based extraction function takes stage-specific input context $\mathcal{X}_t$ and a prompt $\mathcal{P}_t$, and produces $\mathcal{K}_t$:
\begin{equation}
\mathcal{K}_t = f_{\text{LLM}}(\mathcal{X}_t, \mathcal{P}_t), \quad t \in \{1,\ldots,6\}
\end{equation}
where $\mathcal{X}_t$ may include the schema $\mathcal{S}$, a sampled subset of instances $\mathcal{I}_{\text{sample}} \subset \mathcal{I}$, and previously extracted knowledge $\mathcal{K}_{<t}$. The complete knowledge base is then given by $\mathcal{K} = \mathcal{K}_{1:6}$.

The inferred $\mathcal{K}$ is retained in its six layers, and each component retrieves only the layers relevant to its task: the synthesizer consults $\mathcal{K}$ when composing each triple (Section~\ref{subsec:generation}), while the verifier retrieves evidence from the layers that match the class of each diagnosis (Section~\ref{subsec:refinement}). Because $\mathcal{K}$ is inferred rather than declared, its accuracy is audited by human annotators in Section~\ref{subsec:kb-audit}.

\begin{table}[!ht]
\centering
\caption{Overview of the Six-Stage Semantic Extraction Process}
\label{tab:extraction-pipeline}
\resizebox{\columnwidth}{!}{%
\begin{tabular}{lll}
\toprule
\textbf{Stage ($t$)} & \textbf{Input Context} & \textbf{Output ($\mathcal{K}_t$)} \\
\midrule
1. Schema Extraction & $\mathcal{S}, \mathcal{I}_{\text{sample}}$ & $\mathcal{K}_1$ (Metadata) \\
2. Domain Analysis & $\mathcal{S}, \mathcal{K}_1$ & $\mathcal{K}_2$ (Domain Constraints) \\
3. Field Type Analysis & $\mathcal{K}_{1}, \mathcal{I}_{\text{sample}}$ & $\mathcal{K}_3$ (Field Types) \\
4. Column Analysis & $\mathcal{K}_{1:3}$ & $\mathcal{K}_4$ (Column Semantics) \\
5. Table Analysis & $\mathcal{K}_{1:4}$ & $\mathcal{K}_5$ (Table Constraints) \\
6. Relation Analysis & $\mathcal{K}_{3}, \mathcal{K}_{4}, \mathcal{K}_{5}$ & $\mathcal{K}_6$ (Cross-Table Relations) \\
\bottomrule
\end{tabular}%
}
\end{table}

\subsection{Generating Trace-Grounded Drafts}
\label{subsec:generation}

Guided by the knowledge base $\mathcal{K}$, the synthesizer composes each draft as a triple $(q, s, r)$. Its design is shaped by two requirements: the corpus must cover diverse question types and SQL structures, and every draft must stay within what the domain supports. The synthesizer meets the former with generation controls, which decide what to ask, and the latter with a two-step composition order, which decides how the query is derived.

\begin{table}[!ht]
\centering
\caption{Generation control parameters.}
\label{tab:generation-config}
\footnotesize
\setlength{\tabcolsep}{8pt}
\begin{tabular}{ll}
\toprule
\textbf{Parameter} & \textbf{Description} \\
\midrule
$\mathcal{C}$ & Domain context (e.g., sales, education) \\
$\mathcal{T}$ & Analytical task type (e.g., trend analysis, ranking) \\
$\mathcal{L}$ & SQL complexity level ($L_1$--$L_4$; Table~\ref{tab:complexity-levels}) \\
\bottomrule
\end{tabular}
\end{table}

Each draft is instantiated from a generation scenario $\sigma = (\mathcal{C}, \mathcal{T})$ and a complexity level $\ell \in \mathcal{L}$: $\mathcal{C}$ is the domain context (e.g., sales, education), $\mathcal{T}$ is the analytical task type (e.g., trend analysis, ranking), and $\ell$ ranges from single-table queries ($L_1$) to modular compositions with CTEs ($L_4$) (Tables~\ref{tab:generation-config} and~\ref{tab:complexity-levels}). Each draft uses a newly sampled $(\sigma, \ell)$, so the corpus spreads across contexts, task types, and difficulty levels rather than concentrating on a few frequent patterns.

\begin{table}[!ht]
\centering
\caption{SQL Complexity Levels}
\label{tab:complexity-levels}
\footnotesize
\begin{tabular}{ll}
\toprule
\textbf{Level} & \textbf{Query Features} \\
\midrule
$L_1$ & Single-table queries with filtering and aggregation \\
$L_2$ & Multi-table joins \\
$L_3$ & Nested subqueries and window functions \\
$L_4$ & CTEs and modular composition \\
\bottomrule
\end{tabular}
\end{table}

The question and its rationale are composed first. Composing the question against $\mathcal{K}$ keeps it within what the domain supports, and the rationale records the operation the question intends. For the sampled $(\sigma, \ell)$, the prompt incorporates relevant constraints from $\mathcal{K}$, and the synthesizer composes:
\begin{equation}
(q, r) = f_{\text{que}}(\sigma, \ell, \mathcal{K})
\end{equation}
where $q$ is the natural language question and $r$ is the rationale. The constraints from $\mathcal{K}$ serve two purposes: (1) specifying valid field operations (e.g., aggregation-compatible columns), and (2) enforcing domain rules (e.g., excluding refunded transactions from revenue).

Given $(q, r)$ and the same $(\sigma, \ell)$, the SQL query is then derived from the stated intent:
\begin{equation}
s = f_{\text{sql}}(q, r, \sigma, \ell, \mathcal{K})
\end{equation}
Conditioning $s$ on the rationale ties the query to the intent of $q$ and leaves a reasoning record for later inspection. Although composed under the guidance of $\mathcal{K}$, the drafts may still contain semantic errors, and every draft is passed to the verifier (Section~\ref{subsec:refinement}).

\subsection{Diagnosing and Refining Semantic Errors}
\label{subsec:refinement}

The verifier decides which drafts enter the corpus. Its objective is to separate the failures the database reports by itself from the semantic errors that only the question and the domain knowledge can reveal, and to repair what is repairable. This section defines semantic errors, then details the iterative loop that detects and corrects them.

The checks draw on two kinds of evidence. Execution and schema checks need no additional knowledge: the database reports every failure itself, whether a query does not parse, references a column that does not exist, or violates a constraint declared in the schema. Semantic checks have no such source: \texttt{AVG(CDSCode)} executes without error while averaging a school identifier is meaningless, and this kind of error leaves no trace in execution results~\cite{shen2025study}; it becomes visible only when the question is read against the domain knowledge in $\mathcal{K}$. Data synthesis adds a failure mode that conventional text-to-SQL never faces: the question itself may ask something the domain does not support. Definition~1 therefore applies to both the question and the query.

\textbf{Definition 1 (Semantic Error).} Let $(q,s)$ be a synthesized pair whose SQL query parses, executes, and conforms to the constraints declared in the schema. The pair contains a semantic error if any of the following conditions is violated: (i) \emph{question validity}: the question asks something the domain supports, and its concepts and operations agree with the meanings of the involved columns and with domain rules; (ii) \emph{intent match}: $s$ faithfully expresses the intended operation of $q$, covering its task, conditions, aggregation, granularity, and result form; and (iii) \emph{usage consistency}: $s$ respects the domain knowledge of the database in how it uses each column, value, and join, and violations include averaging an identifier column and joining two unrelated tables.

Definition~1 states what a valid pair must satisfy; Table~\ref{tab:semantic-taxonomy} states how each violation is recognized and repaired. The seven classes partition the two repairable conditions: classes I-1 to I-4 refine intent match into four ways a query can miss the request of the question; classes II-1 and II-2 catch element misuse, where a column is used against what it actually means or a join expresses a relationship that does not exist; and class III-1 catches violations of the rules of the domain. Every check reads the pair $(q,s)$ itself; what distinguishes the classes is the additional evidence required to establish a violation: none for the intent-mismatch classes, the meanings of the involved elements for element misuse, and the rules themselves for a domain-rule violation. Question validity has no dedicated class: an unsupported request, such as asking for the average of an identifier, surfaces as an element-misuse or domain-rule violation attributed to $q$; since the question is never edited, such pairs are rejected rather than repaired. This organization serves correction directly: a diagnosis names the class, the class names the evidence, and the evidence conditions the edit. Full class descriptions and the diagnostic and correction prompts are given in Appendix~\ref{app:taxonomy}.

\begin{table}[!t]
\centering
\caption{Semantic error taxonomy.}
\label{tab:semantic-taxonomy}
\resizebox{\columnwidth}{!}{%
\begin{tabular}{lll}
\toprule
\textbf{ID} & \textbf{Class} & \textbf{Additional evidence} \\
\midrule
I-1 & Task mismatch & none \\
I-2 & Condition mismatch & none \\
I-3 & Aggregation/granularity mismatch & none \\
I-4 & Ordering/result-form mismatch & none \\
\midrule
II-1 & Field-role misuse & $\mathcal{K}_3$, $\mathcal{K}_4$ \\
II-2 & Join misuse & $\mathcal{K}_6$ \\
\midrule
III-1 & Domain-rule violation & $\mathcal{K}_2$, $\mathcal{K}_5$, $\mathcal{K}_6$ \\
\bottomrule
\end{tabular}%
}
\end{table}

Detection and correction proceed in an iterative \textsc{Diagnose}--\textsc{Retrieve}--\textsc{Correct} loop that runs for at most three iterations. Writing the current triple as $\tau^{(k)}=(q^{(k)},s^{(k)},r^{(k)})$ and taking the composed draft as $\tau^{(0)}$, one refinement iteration is formalized as:
\begin{equation}
\label{eq:refinement}
\begin{aligned}
(\mathcal{B}^{(k)},\mathcal{E}^{(k)}) &= \textsc{Diagnose}(\tau^{(k)},\mathcal{K},\mathcal{D}), \\
\Phi^{(k)} &= \textsc{Retrieve}(\mathcal{E}^{(k)},\mathcal{K}), \\
\tau^{(k+1)} &= \textsc{Correct}(\tau^{(k)},\mathcal{B}^{(k)},\Phi^{(k)}),
\end{aligned}
\end{equation}
where $\mathcal{B}^{(k)}$ records the execution and schema errors and $\mathcal{E}^{(k)}$ records the semantic errors.

The \textsc{Diagnose} function takes the current triple $\tau^{(k)}$ together with $\mathcal{K}$ and $\mathcal{D}$, and produces the two lists. The execution and schema checks parse and execute $s^{(k)}$ and verify its references against the schema; every failure is recorded in $\mathcal{B}^{(k)}$. The semantic checks examine the conditions of Definition~1 in their stated order, question validity first: an invalid question leaves no standard for the remaining checks, and the examination stops there. The $m_k$ confirmed violations are each labeled with one of the seven classes of Table~\ref{tab:semantic-taxonomy}, collected in $\mathcal{E}^{(k)}$, and written as $e_j^{(k)}=(t_j,u_j,\lambda_j,v_j)$, which records the class, the affected artifact, the location, and the violated condition; one pair may carry several violations at once.

The \textsc{Retrieve} function fetches, for each diagnosis $e_j^{(k)}$, the evidence needed to correct it. The class $t_j$ names the additional evidence to draw, as listed in Table~\ref{tab:semantic-taxonomy}: element misuse draws on what the involved elements mean ($\mathcal{K}_3$, $\mathcal{K}_4$, $\mathcal{K}_6$), and a domain-rule violation draws on the rules themselves ($\mathcal{K}_2$, $\mathcal{K}_5$, $\mathcal{K}_6$). Intent-mismatch diagnoses are established from the pair alone; knowledge entries are added only when reading the question or the implicated clause involves element or domain semantics. Because $\mathcal{K}$ is inferred rather than declared, treating it as evidence is an operational assumption; its support is audited in Section~\ref{subsec:kb-audit}. The evidence is paired with its diagnosis as
\begin{equation}
\label{eq:evidence-retrieval}
\Phi^{(k)}=\left\{\left(e_j^{(k)},\phi_j^{(k)}\right)\right\}_{j=1}^{m_k},
\end{equation}
and the pairs, together with the failures in $\mathcal{B}^{(k)}$, are assembled into the correction prompt.

The \textsc{Correct} function applies local, diagnosis-conditioned edits at the locations the diagnoses point to; it never regenerates the whole triple. The question is never edited: every check, intent match included, is defined relative to $q$, and editing $q$ would dissolve the standard it is judged against. A pair whose question is invalid is therefore unrepairable and is rejected. Within $s$ and $r$, each edit lands where its class points: an intent mismatch is repaired in the clause that implements the mismatched aspect of the question, a condition mismatch in the \texttt{WHERE} or \texttt{HAVING} clause, a granularity mismatch in the aggregate expression and \texttt{GROUP BY}; an element misuse or domain-rule violation is repaired at the implicated column, value, or join. The failures in $\mathcal{B}^{(k)}$ are corrected from the database feedback that reported them. After each edit, the rationale is updated to stay consistent with the corrected query, and the applied edits are retained in a refinement record.

After each correction, all checks are reapplied. A triple enters the synthetic corpus only when every check passes; a triple with unresolved violations at the iteration limit is likewise rejected rather than counted as corrected.

\section{Experiments}
\label{sec:experiments}

Our empirical study answers four research questions. \textbf{RQ1 (downstream performance):} do models fine-tuned on SemanticAgent data outperform those trained on existing synthetic corpora? \textbf{RQ2 (final corpus quality):} does the final synthesized corpus contain fewer semantic errors, and what execution-level and distributional properties does it exhibit? \textbf{RQ3 (mechanism and reliability):} can the knowledge base, diagnosis and repair process, and external controls account for the observed gains? \textbf{RQ4 (scaling):} how does downstream performance change with the synthetic data budget? We finally reanalyze the anomalous Semantic Alignment result on BIRD.

\subsection{Experimental Setup}

\subsubsection{Datasets and Baselines}

To evaluate SemanticAgent, we conduct experiments on eight benchmarks: Spider, BIRD, and Spider~2.0 as the main benchmarks, and EHRSQL, ScienceBenchmark, Spider-Syn, Spider-Realistic, and Spider-DK probing domain transfer, question robustness, and domain knowledge.

Spider~\cite{yu2018spider} is the standard cross-domain benchmark, with 10,181 question--SQL pairs over 200 databases from 138 domains. BIRD~\cite{li2023can} contains 12,751 question--SQL pairs over 95 databases in 37 domains; 42.8\% of its questions require knowledge beyond direct schema matching, probing knowledge-intensive reasoning. Spider~2.0~\cite{lei2024spider2} targets enterprise-level scenarios with large schemas and multi-step SQL workflows; following prior work, we use its SQLite subset (135 of 632 tasks).

EHRSQL~\cite{lee2022ehrsql} and ScienceBenchmark~\cite{zhang2023sciencebenchmark} assess transfer to the clinical and scientific domains; Spider-Syn~\cite{gan2021syn} and Spider-Realistic~\cite{deng2021structure} test robustness against synonym substitution and explicit column-name removal; and Spider-DK~\cite{gan2021exploring} targets questions that require domain knowledge beyond schema-level matching.

SemanticAgent is compared against three synthetic-data generation methods under a controlled protocol: CodeS~\cite{li2024codes}, a template-based method that increases structural diversity across schemas; SynQL~\cite{baumgartner2024synql}, which synthesizes SQL by decomposing and recombining abstract syntax trees; and OmniSQL~\cite{li2025omnisql}, which combines structural transformations with LLM-based verification. All methods synthesize from the same source databases under an identical data budget, and the resulting corpora are fine-tuned and evaluated with the same downstream backbone and training configuration, so corpus quality is the only intended difference.

Two correction-capable frameworks are additionally adapted to this comparison. MAC-SQL~\cite{wang2025macsql} and MAGIC~\cite{askari2025magic} take a question as input and generate or correct SQL at inference time; fed with the same question set, they produce comparable corpora, which are evaluated under the same protocol (Section~\ref{subsec:generator-control}). The released OmniSQL models and prompting-based systems (DIN-SQL~\cite{pourreza2023din}, CHESS~\cite{talaei2024chess}) are included as reference points rather than controlled baselines.

\subsubsection{Synthetic Corpora and Contamination Control}

\begin{table}[!t]
\centering
\caption{Database Statistics and Synthetic Corpus Sizes}
\label{tab:synth_source_stats}
\resizebox{\columnwidth}{!}{%
\begin{tabular}{lccccc}
\toprule
Benchmark & DB & Domain & Avg. Tbl/DB & Avg. Col/DB & Synth. Pairs \\
\midrule
Spider 1.0 & 200 & 138 & 5.11 & 26.82 & 40,000 \\
Spider~2.0-SQLite & 30 & 8 & 13.8 & 7.0 & 20,000 \\
BIRD & 95 & 37 & 7.64 & 54.56 & 30,000 \\
EHRSQL & 2 & 1 & 31.55 & 92.00 & 5,000 \\
ScienceBenchmark & 3 & 3 & 16.67 & 44.47 & 5,000 \\
\bottomrule
\end{tabular}%
}
\end{table}

Synthetic corpora are constructed separately for each benchmark, scaled to its size and schema complexity; Table~\ref{tab:synth_source_stats} reports the number of source databases and domains, the average tables and columns per database, and the resulting corpus sizes. For Spider~2.0-SQLite, which provides only evaluation databases without a training split, we synthesize directly from its SQLite databases; Spider-Syn, Spider-Realistic, and Spider-DK are reserved for evaluation only.

All pairs are generated exclusively from database-side information, the schema and sampled cell values; no gold question--SQL pairs from any public benchmark are used, neither as synthesis seeds nor as in-context exemplars. Synthesizing from the databases used for evaluation therefore does not by itself constitute leakage: a schema affords a large space of distinct question--SQL pairs, and the generated pairs do not replicate gold annotations.

As a second safeguard against accidental overlap, we discard any synthesized pair whose question or SQL shares an n-gram overlap above [THRESHOLD] with an evaluation instance. This filter removes 764 pairs (0.76\% of the initially generated corpus), leaving 99,236 pairs for downstream training.

\subsubsection{Evaluation Metrics and Human Evaluation}
\label{subsec:human-evaluation}\label{subsec:direct-semantic-metrics}

\paragraph{Overview.}
For RQ1, we report execution accuracy (EX)~\cite{yu2018spider,li2023can} on the remaining benchmarks and test-suite accuracy (TS)~\cite{zhong2020semantic} on Spider (dev), Spider-Syn, and Spider-Realistic. For RQ2, the primary readouts are the verifiable semantic error rate and the frequencies of the seven semantic-error classes from direct human evaluation; SER, SQL complexity, and embedding-based diversity provide complementary corpus statistics. Semantic Alignment (SA) is an auxiliary measure and is reanalyzed on BIRD in Section~\ref{subsec:bird-sa-reanalysis}. RQ3 reuses the semantic error rate and downstream EX under controlled configurations, and RQ4 analyzes EX as a function of the synthetic data budget. Intrinsic-evaluation prompts are provided in Appendix~\ref{app:prompts}.

\paragraph{Protocol.}
The human evaluation covers executable pairs sampled from the final corpora of CodeS, SynQL, OmniSQL, and SemanticAgent. Annotators receive the question, SQL query, source database, and independently traceable domain evidence when needed; the automatically inferred knowledge base and system diagnoses are not treated as ground truth. Each pair is judged under Definition~1 and labeled with all applicable classes of Table~\ref{tab:semantic-taxonomy}; insufficiently supported cases are marked not assessable. For the paired refinement analysis, initial and final versions of the same SemanticAgent drafts are labeled under the same criteria. Two annotators label independently after a pilot alignment, and a third expert adjudicates disagreements (inter-annotator agreement [K]). The full guideline is given in Appendix~\ref{app:annotation-guideline}.

\paragraph{Metrics.}
With $\mathcal{H}_{\mathrm{assess}}$ denoting the sampled pairs judged answerable, executable, and sufficiently supported by independent evidence, and $\operatorname{Err}(x)$ denoting the set of semantic-error labels assigned to pair $x$, the verifiable semantic error rate is:

\begin{equation}
\label{eq:verifiable-semantic-error-rate}
\operatorname{ErrRate}_{\mathrm{sem}} = \frac{\left|\left\{x\in\mathcal{H}_{\mathrm{assess}} : \operatorname{Err}(x)\neq\varnothing\right\}\right|}{\left|\mathcal{H}_{\mathrm{assess}}\right|}.
\end{equation}

For each error class $t$, we additionally report:

\begin{equation}
\label{eq:per-type-semantic-error-frequency}
\operatorname{Freq}(t) = \frac{\left|\left\{x\in\mathcal{H}_{\mathrm{assess}} : t\in\operatorname{Err}(x)\right\}\right|}{\left|\mathcal{H}_{\mathrm{assess}}\right|}.
\end{equation}

The denominator is defined separately for each evaluated corpus as its assessable sampled pairs. Because the taxonomy is multi-label, class frequencies need not sum to the overall rate. SER is reported separately as an execution-level measure.

\subsubsection{Implementation Details}
\label{subsec:implementation_details}

For synthetic data generation and LLM-based intrinsic evaluation, we use Qwen3-235B-A22B-Instruct~\cite{yang2025qwen3}, deployed locally with vLLM~0.11.0~\cite{kwon2023efficient}. For downstream evaluation, we fully fine-tune three backbone models: Qwen2.5-Coder-7B, Qwen2.5-Coder-14B~\cite{hui2024qwen25coder}, and Qwen3-14B~\cite{yang2025qwen3}. All downstream backbones are trained on each synthetic corpus under identical settings.

We use the ms-swift framework~\cite{zhao2025swift} with DeepSpeed ZeRO Stage~3 offloading~\cite{rajbhandari2021zero} and \texttt{bfloat16} mixed precision. We optimize all downstream models with AdamW~\cite{loshchilov2019decoupled}, using $\beta_1=0.9$, $\beta_2=0.95$, and $\epsilon=10^{-8}$. The learning rate is linearly warmed up over the first $5\%$ of training steps and then annealed with a cosine schedule to $10\%$ of its peak value. Unless otherwise specified, all models are fully fine-tuned for 2 epochs with an effective global batch size of 512 and a peak learning rate of $4\times10^{-6}$. We use a maximum sequence length of 8{,}192 tokens, weight decay of 0.1, and gradient clipping of 1.0.

All local experiments are conducted on a server equipped with 8 NVIDIA A800-SXM4-80GB GPUs and 1 TB of RAM. PyTorch~2.9.1, vLLM~0.11.0 is used for high-throughput inference in data synthesis and intrinsic evaluation. DeepSpeed~0.18.6 is used for distributed training.


\subsection{Downstream Performance (RQ1)}
\label{subsec:downstream_performance}

\subsubsection{Matched Synthesis Comparison}
\label{subsec:matched-comparison}

As shown in Table~\ref{tab:main_results}, the controlled synthesis comparison evaluates corpora generated by CodeS, SynQL, OmniSQL, and SemanticAgent under the same source databases, data budget, downstream backbone, and training settings. The question-controlled comparison with MAC-SQL and MAGIC is reported separately in Section~\ref{subsec:generator-control}, because those frameworks receive a fixed question set and use their own generation and correction procedures.
\begin{table*}[t]
\centering
\caption{Downstream performance on text-to-SQL benchmarks. All values are EX (\%), except the Spider columns, which report TS (\%); higher is better. Bold marks the best result within each block of the controlled synthesis comparison; reference and published rows are shown for context.}
\label{tab:main_results}
\small
\setlength{\tabcolsep}{3pt}
\renewcommand{\arraystretch}{1.05}
\begin{tabular}{@{}lccccccccc@{}}
\toprule
\multirow{2}{*}{\textbf{Method}}
& \multicolumn{2}{c}{\textbf{Spider}}
& \multirow{2}{*}{\textbf{BIRD}}
& \multirow{2}{*}{\textbf{S2-SQLite}}
& \multirow{2}{*}{\textbf{Science}}
& \multirow{2}{*}{\textbf{EHRSQL}}
& \multicolumn{3}{c}{\textbf{Spider Variants}} \\
\cmidrule(lr){2-3} \cmidrule(lr){8-10}
& \textbf{Dev} & \textbf{Test} & & & & & \textbf{DK} & \textbf{Syn} & \textbf{Real.} \\
\midrule
\multicolumn{10}{l}{\emph{Reference systems}} \\
\midrule
QC2.5-7B         & 73.4 & 82.2 & 50.9 & 1.5  & 45.2 & 24.3 & 67.5 & 63.1 & 66.7 \\
QC2.5-14B        & 78.1 & 86.6 & 61.5 & 5.9  & 52.2 & 31.6 & 73.6 & 68.2 & 76.2 \\
Q3-14B           & 76.7 & 83.9 & 55.2 & 1.4  & 53.1 & 38.5 & 74.0 & 67.1 & 71.3 \\
QC2.5-32B        & 77.7 & 87.5 & 64.5 & 5.9  & 54.8 & 36.4 & 78.3 & 69.9 & 72.1 \\
Q2.5-72B         & 73.9 & 84.0 & 60.3 & 9.6  & 52.8 & 35.0 & 76.4 & 64.1 & 70.1 \\
DeepSeek-V3      & 73.1 & 85.5 & 63.2 & 12.6 & 56.2 & 43.2 & 72.9 & 64.4 & 67.9 \\
DIN-SQL (GPT-4o) & 74.2 & 85.3 & 55.9 & 9.1  & 59.7 & 42.5 & 74.5 & 63.8 & 68.2 \\
CHESS (GPT-4o)   & 76.8 & 87.2 & 66.8 & 7.7  & 61.3 & 48.2 & 77.5 & 66.9 & 72.4 \\
\midrule
\multicolumn{10}{l}{\emph{Published specialized models}} \\
\midrule
OmniSQL-7B       & 81.2 & 87.9 & 63.9 & 10.4 & 50.2 & 34.9 & 76.1 & 69.7 & 76.2 \\
OmniSQL-14B      & 81.4 & 88.3 & 64.2 & 10.4 & 56.9 & 39.9 & 72.9 & 69.0 & 76.4 \\
OmniSQL-32B      & 80.9 & 87.6 & 64.5 & 11.9 & 57.2 & 42.4 & 76.1 & 69.7 & 78.1 \\
\midrule
\multicolumn{10}{l}{\emph{Controlled synthesis comparison}} \\
\midrule
CodeS (QC2.5-7B)         & 76.1 & 82.8 & 57.2 & 2.1  & 46.8 & 30.2 & 70.3 & 64.9 & 69.8 \\
SynQL (QC2.5-7B)         & 77.9 & 83.0 & 59.6 & 3.7  & 48.5 & 33.1 & 72.0 & 66.4 & 72.3 \\
OmniSQL (QC2.5-7B)       & 76.1 & 84.9 & 59.8 & 7.1  & 47.4 & 34.5 & 72.6 & 64.8 & \textbf{72.7} \\
MAC-SQL (QC2.5-7B)       & $[\cdot]$ & $[\cdot]$ & $[\cdot]$ & $[\cdot]$ & $[\cdot]$ & $[\cdot]$ & $[\cdot]$ & $[\cdot]$ & $[\cdot]$ \\
MAGIC (QC2.5-7B)         & $[\cdot]$ & $[\cdot]$ & $[\cdot]$ & $[\cdot]$ & $[\cdot]$ & $[\cdot]$ & $[\cdot]$ & $[\cdot]$ & $[\cdot]$ \\
\textbf{Ours (QC2.5-7B)} & 76.3 & \textbf{86.1} & \textbf{62.4} & \textbf{9.3} & \textbf{50.6} & \textbf{36.8} & \textbf{73.9} & \textbf{67.5} & 72.4 \\
\cmidrule{1-10}
CodeS (QC2.5-14B)        & 79.5 & 87.3 & 62.1 & 4.3  & 53.4 & 34.0 & 72.1 & 69.0 & 77.0 \\
SynQL (QC2.5-14B)        & 80.6 & 87.8 & 63.2 & 5.8  & 54.6 & 36.1 & 74.8 & 69.8 & \textbf{78.0} \\
OmniSQL (QC2.5-14B)      & \textbf{81.2} & \textbf{88.2} & 64.0 & 10.2 & 56.4 & 38.8 & 73.2 & \textbf{70.2} & 75.5 \\
MAC-SQL (QC2.5-14B)      & $[\cdot]$ & $[\cdot]$ & $[\cdot]$ & $[\cdot]$ & $[\cdot]$ & $[\cdot]$ & $[\cdot]$ & $[\cdot]$ & $[\cdot]$ \\
MAGIC (QC2.5-14B)        & $[\cdot]$ & $[\cdot]$ & $[\cdot]$ & $[\cdot]$ & $[\cdot]$ & $[\cdot]$ & $[\cdot]$ & $[\cdot]$ & $[\cdot]$ \\
\textbf{Ours (QC2.5-14B)}& 80.1 & 86.3 & \textbf{65.4} & \textbf{10.8} & \textbf{57.8} & \textbf{40.6} & \textbf{76.5} & 70.1 & 75.2 \\
\cmidrule{1-10}
OmniSQL (Q3-14B)         & \textbf{81.1} & 86.8 & 63.7 & 9.6  & 52.6 & 37.8 & 75.1 & 67.6 & \textbf{74.8} \\
\textbf{Ours (Q3-14B)}   & 81.0 & \textbf{88.1} & \textbf{66.0} & \textbf{13.0} & \textbf{56.5} & \textbf{40.1} & \textbf{76.8} & \textbf{69.2} & 74.0 \\
\bottomrule
\end{tabular}

\vspace{1mm}
\begin{minipage}{\textwidth}
\footnotesize
\textit{Note:} \emph{Reference systems} (backbones and prompting-based systems) and \emph{published specialized models} (fine-tuned on their own synthetic data in prior work) are shown for context and are not under the matched protocol. The \emph{controlled synthesis comparison} builds every corpus from the same source databases under an identical data budget, then fine-tunes and evaluates all corpora identically; the MAC-SQL and MAGIC rows follow the question-controlled protocol of Section~\ref{subsec:generator-control}, where these frameworks generate SQL for the same question set. QC2.5-xB = Qwen2.5-Coder-xB; Q2.5-72B = Qwen2.5-72B; Q3-14B = Qwen3-14B; S2-SQLite = Spider~2.0-SQLite.
\end{minipage}

\vspace{-2mm}
\end{table*}


The gains over OmniSQL, the strongest prior synthesis method, concentrate on the most challenging benchmarks. On BIRD, SemanticAgent improves over OmniSQL by 2.6, 1.4, and 2.3 points under Qwen2.5-Coder-7B, Qwen2.5-Coder-14B, and Qwen3-14B, respectively. On Spider~2.0-SQLite, the corresponding gains are 2.2, 0.6, and 3.4 points. These results suggest that SemanticAgent is particularly effective when schema complexity and compositional reasoning demands are high.

On Spider, the margin is smaller. SemanticAgent improves over OmniSQL in the Qwen2.5-Coder-7B and Qwen3-14B settings (86.1 vs.\ 84.9 and 88.1 vs.\ 86.8), but is slightly below OmniSQL with Qwen2.5-Coder-14B (86.3 vs.\ 88.2). This weaker separation is not surprising, given Spider's maturity and the relatively limited headroom after synthetic fine-tuning.

The benefits extend to specialized domains. On ScienceBenchmark, SemanticAgent surpasses OmniSQL by 1.4--3.9 points across the three backbone settings, and on EHRSQL the corresponding gains range from 1.8 to 2.3 points, suggesting that the synthesized data better capture domain-specific semantics.

A similar pattern appears on the robustness variants. On Spider-Syn, SemanticAgent improves over OmniSQL by 2.7 points under Qwen2.5-Coder-7B and 1.6 points under Qwen3-14B; on Spider-DK, it improves in all three settings, by 1.3, 3.3, and 1.7 points. On Spider-Realistic, it remains competitive but does not consistently achieve the best score, suggesting that synthetic supervision improves robustness to lexical variation more clearly than robustness to natural user phrasing.

The two correction-capable generators are evaluated separately under the fixed-question protocol of Section~\ref{subsec:generator-control}; their semantic-quality and downstream results are reported there.

Overall, the matched comparison supports two conclusions. First, under identical budgets and protocols, SemanticAgent improves over the prior synthesis methods included in this comparison. Second, the gains concentrate on benchmarks that demand stronger schema understanding, domain-specific reasoning, or robustness to lexical variation, indicating that higher-quality synthetic supervision is most valuable in semantically demanding settings.

\subsubsection{Reference Systems and Published Models}

Beyond the matched protocol, Table~\ref{tab:main_results} reports two context groups. The released OmniSQL models, fine-tuned by their authors on their own corpora, reach 80.9--81.4 on Spider (dev) and 63.9--64.5 on BIRD; the gap between these rows and the controlled OmniSQL rows likely reflects the training budget and protocol rather than the synthesis method, which is why they are not treated as direct baselines. SemanticAgent also remains competitive with much larger general-purpose or prompting-based systems: the Qwen2.5-Coder-7B variant reaches 62.4 on BIRD, outperforming Qwen2.5-72B~\cite{yang2024qwen25} (60.3), and the Qwen3-14B variant achieves 13.0 on Spider~2.0, surpassing DeepSeek-V3~\cite{liu2024deepseekv3} (12.6), DIN-SQL~\cite{pourreza2023din} (9.1), and CHESS~\cite{talaei2024chess} (7.7). Because these comparisons are not controlled, they are read as reference points only.

\subsection{Data Quality (RQ2)}

\subsubsection{Direct Evaluation of Semantic Errors}
\label{subsec:direct-semantic-evaluation}

We directly evaluate the semantic quality of the final synthesized corpora through the human-evaluation protocol of Section~\ref{subsec:human-evaluation}. The evaluation covers executable pairs sampled from the final corpora of CodeS, SynQL, OmniSQL, and SemanticAgent. Table~\ref{tab:verifiable-semantic-error-rate} reports the overall semantic error rate, and Table~\ref{tab:semantic-error-type-frequency} reports the frequency of each error class.

\begin{table}[!t]
\centering
\caption{Human-evaluated semantic error rates in final synthesized corpora. Error rates are computed over assessable sampled pairs. Values marked with $[\cdot]$ are placeholders to be replaced by the measured results.}
\label{tab:verifiable-semantic-error-rate}
\small
\setlength{\tabcolsep}{4pt}
\resizebox{\columnwidth}{!}{%
\begin{tabular}{lrrrr}
\toprule
\textbf{Corpus} & \textbf{Evaluated} & \textbf{Assessable} & \textbf{Errors} & \textbf{Error rate (\%)} \\
\midrule
CodeS & [N] & [N] & [N] & [P] \\
SynQL & [N] & [N] & [N] & [P] \\
OmniSQL & [N] & [N] & [N] & [P] \\
SemanticAgent & [N] & [N] & [N] & [P] \\
\bottomrule
\end{tabular}%
}
\end{table}

Table~\ref{tab:verifiable-semantic-error-rate} shows that [P1]\% of the sampled executable CodeS pairs, [P2]\% of the SynQL pairs, and [P3]\% of the OmniSQL pairs violate Definition~1. The final SemanticAgent corpus has an error rate of [P4]\%, compared with [P3]\% for the strongest baseline, corresponding to a relative reduction of [P-REL]\%.

\begin{table*}[!t]
\centering
\caption{Per-class frequencies of verifiable semantic errors in final synthesized corpora. Frequencies are computed over assessable sampled pairs. Error classes are non-exclusive, so their frequencies do not sum to the overall error rate. Values marked with $[\cdot]$ are placeholders to be replaced by the measured results.}
\label{tab:semantic-error-type-frequency}
\small
\setlength{\tabcolsep}{4pt}
\begin{tabular}{lrrrrrrr}
\toprule
\textbf{Corpus} & \textbf{I-1} & \textbf{I-2} & \textbf{I-3} & \textbf{I-4} & \textbf{II-1} & \textbf{II-2} & \textbf{III-1} \\
\midrule
CodeS & [n (\%)] & [n (\%)] & [n (\%)] & [n (\%)] & [n (\%)] & [n (\%)] & [n (\%)] \\
SynQL & [n (\%)] & [n (\%)] & [n (\%)] & [n (\%)] & [n (\%)] & [n (\%)] & [n (\%)] \\
OmniSQL & [n (\%)] & [n (\%)] & [n (\%)] & [n (\%)] & [n (\%)] & [n (\%)] & [n (\%)] \\
SemanticAgent & [n (\%)] & [n (\%)] & [n (\%)] & [n (\%)] & [n (\%)] & [n (\%)] & [n (\%)] \\
\bottomrule
\end{tabular}
\end{table*}

Table~\ref{tab:semantic-error-type-frequency} identifies the error classes that account for the corpus-level differences. The reduction is concentrated in [ERROR CLASSES], while [UNIMPROVED OR REVERSED CLASSES]. The diagnosis and repair analysis in Section~\ref{subsec:diagnosis-evaluation} examines how these labels change during refinement, and the component ablation in Section~\ref{subsec:ablation} tests the contribution of knowledge conditioning.

\subsubsection{Execution-Level and Distributional Properties}

This subsection reports corpus properties that complement the direct human evaluation of Section~\ref{subsec:direct-semantic-evaluation}. SER measures executability, SQL complexity describes structural coverage, and embedding-based diversity describes the spread of the synthesized query space. SA is excluded from these summaries and reanalyzed separately in Section~\ref{subsec:bird-sa-reanalysis}.

\begin{table}[!ht]
\centering
\caption{Successful execution rate (SER) of synthesized corpora. Higher values indicate a larger proportion of executable synthesized SQL queries.}
\label{tab:consistency_metrics}
\resizebox{\columnwidth}{!}{%
\begin{tabular}{lccc}
\toprule
\textbf{Method} & \textbf{BIRD} & \textbf{Spider 1.0} & \textbf{Spider~2.0-SQLite} \\
\midrule
CodeS & 91.37 & 92.73 & 21.06 \\
SynQL & 90.12 & \textbf{94.18} & 25.78 \\
OmniSQL & \textbf{94.83} & \underline{93.28} & \underline{35.31} \\
SemanticAgent & \underline{92.14} & 93.18 & \textbf{63.13} \\
\bottomrule
\end{tabular}%
}
\begin{minipage}{0.5\textwidth}
\footnotesize
\vspace{1mm}
\textit{Note:} SER measures executability only and is not a substitute for the human-evaluated semantic error rate.
\end{minipage}
\end{table}

Table~\ref{tab:consistency_metrics} shows that SER varies across benchmarks rather than uniformly favoring one method. SemanticAgent achieves the highest SER on Spider~2.0-SQLite, where execution is most difficult, but not on every benchmark. This execution-level result is interpreted separately from semantic quality.

\begin{figure*}[!t]
    \centering
    \includegraphics[width=0.62\textwidth]{data_quality.pdf}
    \caption{Structural complexity distribution of synthesized SQL queries across synthesis methods. The distribution uses the four AST-based complexity levels defined in Section~\ref{subsec:generation}.}
    \label{fig:data_quality}
\end{figure*}

Figure~\ref{fig:data_quality} shows that SemanticAgent generates a larger proportion of moderately complex and complex SQL queries, indicating broader structural coverage. Table~\ref{tab:diversity_metrics} reports the distributional variation of the synthesized queries in the embedding space.

\begin{table}[ht]
\centering
\caption{Embedding-based diversity of synthesized corpora. Higher values indicate greater diversity under the corresponding measure.}
\label{tab:diversity_metrics}
\resizebox{\columnwidth}{!}{%
\begin{tabular}{@{}lcccccc@{}}
\toprule
\textbf{Method} & \multicolumn{2}{c}{\textbf{BIRD}} & \multicolumn{2}{c}{\textbf{Spider 1.0}} & \multicolumn{2}{c}{\textbf{Spider~2.0-SQLite}} \\
\cmidrule(lr){2-3} \cmidrule(lr){4-5} \cmidrule(lr){6-7}
& \textbf{Mean L2} & \textbf{1-NN} & \textbf{Mean L2} & \textbf{1-NN} & \textbf{Mean L2} & \textbf{1-NN} \\
\midrule
CodeS & 0.6213 & 0.2441 & 0.6152 & 0.2394 & 0.4086 & 0.2359 \\
SynQL & 0.6421 & 0.2553 & 0.6354 & 0.2519 & 0.5289 & 0.2175 \\
OmniSQL & \underline{0.6623} & \underline{0.2681} & \underline{0.6568} & \underline{0.2421} & \underline{0.5497} & \underline{0.2210} \\
SemanticAgent & \textbf{0.6819} & \textbf{0.2745} & \textbf{0.6737} & \textbf{0.2645} & \textbf{0.6634} & \textbf{0.2792} \\
\bottomrule
\end{tabular}%
}
\begin{minipage}{0.5\textwidth}
\footnotesize
\vspace{1mm}
\textit{Note:} Higher values indicate greater diversity under the corresponding measure.
\end{minipage}
\end{table}

SemanticAgent achieves the highest Mean L2 and 1-NN values on all three benchmarks. Together with the complexity distribution, these results indicate broader structural and distributional coverage. They describe corpus properties only; the direct human evaluation remains the primary evidence for semantic quality.

\subsection{Mechanism and Reliability Analyses (RQ3)}

\subsubsection{Knowledge Base Audit}
\label{subsec:kb-audit}

Everything downstream of construction draws on the knowledge base $\mathcal{K}$: drafts are generated under its guidance, and errors are diagnosed and corrected with its evidence. $\mathcal{K}$ is constructed automatically from the schema and sampled instances, and apart from the schema metadata ($\mathcal{K}_1$) every layer is LLM-inferred and not guaranteed to be correct item by item. We audit these inferred layers directly, entry by entry.

\textbf{Entry-level audit.} We sample [N-AUDIT] entries per knowledge layer from each database. Only $\mathcal{K}_1$ is taken directly from the schema and verified programmatically; every other layer is produced by LLM analysis and is audited by humans. Human experts judge each entry against the database itself, not against the model's own statements. A $\mathcal{K}_2$ entry describes what the database is for---its business domain and terminology---and is checked against the schema as a whole and the official documentation. A $\mathcal{K}_3$ entry claims the semantic type of a field and is checked by inspecting the stored values of the column. A $\mathcal{K}_4$ entry describes what a column stores and is checked against the values, the schema comment, and the documentation. A $\mathcal{K}_5$ entry claims a table-level constraint and is checked by reasoning over the data: the claim is compiled into a violation query, e.g., a duplicate check for a uniqueness claim. A $\mathcal{K}_6$ entry claims how two tables relate; every database has one underlying entity--relation structure, and the claim is checked against it, via the declared keys, or join tests where keys are undeclared.

Two annotators independently check every sampled entry; a third expert adjudicates disagreements. Each entry receives one of four labels: \emph{correct and supported}, \emph{incorrect}, \emph{plausible but unsupported}, or \emph{not assessable}. A contradicting check yields \emph{incorrect}; a claim that no check settles yields \emph{plausible but unsupported}, which is never counted as correct; a check that cannot be run, e.g., an empty column, yields \emph{not assessable}. The audit judges only the entries that were generated; knowledge the database has but extraction missed is not counted. Table~\ref{tab:kb-audit} reports the label distribution per layer, with Wilson 95\% confidence intervals for the correct-and-supported proportion.

% [发现段模板（FActScore 式，有结果后照写，一段一个发现）：
%  1. \textbf{Most entries are verifiable.} Correct-and-supported
%     ranges from [X]\% to [Y]\% across layers, indicating that ...
%  2. \textbf{[K_X] is the weakest layer.} [Z]\% [L, U] of its entries
%     are ..., we therefore narrow ... （措辞：CI 下界过线才说达标）]

\begin{table}[!t]
\centering
\caption{Entry-level audit of the LLM-inferred knowledge base. Values marked with $[\cdot]$ are placeholders to be replaced by the measured results.}
\label{tab:kb-audit}
\small
\setlength{\tabcolsep}{3pt}
\renewcommand{\arraystretch}{1.15}
\resizebox{\columnwidth}{!}{%
\begin{tabular}{lccccc}
\toprule
\textbf{Layer} & \textbf{Sampled} & \textbf{Correct and supported} & \textbf{Incorrect} & \textbf{Plausible but unsupported} & \textbf{Not assessable} \\
\midrule
$\mathcal{K}_2$ Domain constraints & [N] & [n (\%) [L, U]] & [n (\%)] & [n (\%)] & [n (\%)] \\
$\mathcal{K}_3$ Field types & [N] & [n (\%) [L, U]] & [n (\%)] & [n (\%)] & [n (\%)] \\
$\mathcal{K}_4$ Column semantics & [N] & [n (\%) [L, U]] & [n (\%)] & [n (\%)] & [n (\%)] \\
$\mathcal{K}_5$ Table constraints & [N] & [n (\%) [L, U]] & [n (\%)] & [n (\%)] & [n (\%)] \\
$\mathcal{K}_6$ Cross-table relations & [N] & [n (\%) [L, U]] & [n (\%)] & [n (\%)] & [n (\%)] \\
\bottomrule
\end{tabular}%
}
\end{table}

The audit covers a sample, and the pipeline itself does not detect incorrect entries; it does not certify $\mathcal{K}$ as a whole. The contribution of knowledge conditioning is examined in the component ablation (Section~\ref{subsec:ablation}): [RESULT-ABL-KB]. Transferring to a new database requires re-running the six-stage construction ([COST-KB] LLM calls and [TIME-KB] wall-clock hours per database on average); per-database audit results are given in Appendix~\ref{app:kb-stages}.



\subsubsection{Diagnosis and Repair}
\label{subsec:diagnosis-evaluation}\label{subsec:paired-refinement}

The diagnosis stage uses an LLM to identify semantic errors, so we evaluate its predictions against human labels~\cite{bavaresco2025judgebench}. This analysis covers drafts entering refinement, including drafts that later fail to produce an admissible final pair. Each initial draft is labeled with the seven-class taxonomy of Table~\ref{tab:semantic-taxonomy}; when a final version exists, it is labeled under the same criteria.

\paragraph{Diagnosis reliability.}
We compare the KB-conditioned diagnosis module with a generic critic that uses the same model, decoding settings, taxonomy, output schema, and call structure but omits retrieved knowledge. A draft containing at least one human-confirmed semantic error is positive. Table~\ref{tab:diagnosis-performance} reports Precision, Recall, F1, and the false-positive rate; deterministic execution and schema failures are excluded. Table~\ref{tab:recall-by-type} reports recall for the existing error classes, with I-1--I-4 combined in the main text because they are established from the pair $(q,s)$ itself.

\begin{table}[!t]
\centering
\caption{Diagnosis performance against human-labeled initial drafts. A positive instance contains at least one human-verified semantic error. Values marked with $[\cdot]$ are placeholders.}
\label{tab:diagnosis-performance}
\small
\setlength{\tabcolsep}{4pt}
\begin{tabular}{lrrrr}
\toprule
\textbf{Configuration} & \textbf{Precision} & \textbf{Recall} & \textbf{F1} & \textbf{FPR} \\
\midrule
Generic critic & [P] & [R] & [F1] & [FPR] \\
KB-conditioned diagnosis & [P] & [R] & [F1] & [FPR] \\
\bottomrule
\end{tabular}
\end{table}

\begin{table}[!t]
\centering
\caption{Class-level recall of the two diagnosis configurations. The denominator of each row is the drafts with at least one human-confirmed violation of that class. The complete seven-class breakdown is given in Appendix~\ref{app:diagnosis-by-type}. Values marked with $[\cdot]$ are placeholders.}
\label{tab:recall-by-type}
\small
\setlength{\tabcolsep}{4pt}
\begin{tabular}{lrrr}
\toprule
\textbf{Class} & \textbf{$n$} & \textbf{Generic critic} & \textbf{KB-cond.} \\
\midrule
I-1--I-4 intent mismatches & [N] & [R] & [R] \\
II-1 Field-role misuse & [N] & [R] & [R] \\
II-2 Join misuse & [N] & [R] & [R] \\
III-1 Domain-rule violation & [N] & [R] & [R] \\
\bottomrule
\end{tabular}
\end{table}

The KB-conditioned diagnosis reaches [F1-K] F1, with Recall [R-K]\% and FPR [FPR-K]\% (Table~\ref{tab:diagnosis-performance}). Its recall changes from [R-G]\% to [R-K]\% relative to the generic critic. The class-level results show [DIAGNOSIS-RESULT: state whether the change is concentrated in II-1, II-2, and III-1, and whether the FPR remains acceptable].

\paragraph{Changes during refinement.}
We next compare the taxonomy labels assigned to the initial and final versions of the same drafts. The table reports changes in the existing seven semantic-error labels; it does not introduce additional error types. \emph{No admitted final pair} denotes a pipeline disposition, such as rejection at the iteration limit, rather than a semantic-error class.

\begin{table}[!t]
\centering
\caption{Paired transitions of the seven semantic-error labels after refinement. Initial and final states are derived from Table~\ref{tab:semantic-taxonomy}; \emph{no admitted final pair} is a pipeline disposition rather than an additional error type. Percentages are computed within each initial state. Values marked with $[\cdot]$ are placeholders.}
\label{tab:refinement-transitions}
\small
\setlength{\tabcolsep}{4pt}
\begin{tabular}{llrr}
\toprule
\textbf{Initial state} & \textbf{Final state or disposition} & \textbf{Count} & \textbf{Within-state \%} \\
\midrule
\multirow{3}{*}{At least one taxonomy error}
& No taxonomy error & [N] & [P] \\
& At least one taxonomy error & [N] & [P] \\
& No admitted final pair & [N] & [P] \\
\midrule
\multirow{3}{*}{No taxonomy error}
& No taxonomy error & [N] & [P] \\
& At least one taxonomy error & [N] & [P] \\
& No admitted final pair & [N] & [P] \\
\bottomrule
\end{tabular}
\end{table}

Table~\ref{tab:refinement-transitions} shows that [TRANSITION-RESULT: report the shares of error-bearing drafts that become error-free, retain taxonomy errors, or lack an admitted final pair, and the share of clean drafts in which a taxonomy error is introduced]. The per-class label changes and representative cases are reported in Appendix~\ref{app:refinement-cases}. This analysis evaluates diagnosis reliability and label changes within refinement; the component ablation below tests the causal contribution of knowledge conditioning.

\subsubsection{Component Ablation}
\label{subsec:ablation}

We next isolate the contributions of KB-conditioned generation and refinement. Table~\ref{tab:ablation-ladder} compares an enriched-schema baseline without KB evidence (A), the analyzer and KB-conditioned synthesizer without the Verifier (B), the same synthesizer followed by a generic refinement loop without retrieved knowledge (C), and the full pipeline with KB-conditioned refinement (D). The four configurations use the same schema representation, synthesis prompts, data budget, and downstream training protocol.

\begin{table*}[!t]
\centering
\caption{Ablation of the semantic verification pipeline. Adjacent comparisons isolate KB-conditioned generation (B$-$A), an additional refinement stage without retrieved knowledge (C$-$B), and knowledge conditioning within the same refinement loop (D$-$C). The verifiable semantic error rate is the primary quality measure; SER reports retention, and EX reports downstream execution accuracy. Values marked with $[\cdot]$ are placeholders.}
\label{tab:ablation-ladder}
\small
\setlength{\tabcolsep}{4pt}
\renewcommand{\arraystretch}{1.15}
\begin{tabular}{lcccccc}
\toprule
\textbf{Configuration} & \textbf{KB at generation} & \textbf{Refinement loop} & \textbf{KB in loop} & \textbf{SER} & \textbf{Error rate} & \textbf{EX} \\
\midrule
A. Enriched-schema baseline & no & no & not applicable & [SER-A] & [ERR-A] & [EX-A] \\
B. KB-conditioned generation, w/o Verifier & yes & no & not applicable & [SER-B] & [ERR-B] & [EX-B] \\
C. + generic refinement & yes & yes & no & [SER-C] & [ERR-C] & [EX-C] \\
D. Full SemanticAgent & yes & yes & yes & [SER-D] & [ERR-D] & [EX-D] \\
\bottomrule
\end{tabular}
\end{table*}

The full pipeline achieves an error rate of [ERR-D], compared with [ERR-A], [ERR-B], and [ERR-C] for configurations A--C, respectively (Table~\ref{tab:ablation-ladder}). KB-conditioned generation changes the error rate by [DELTA-BA] points (B$-$A). Adding an otherwise identical refinement loop without retrieved knowledge contributes [DELTA-CB] points (C$-$B), while grounding that same loop in the knowledge base contributes a further [DELTA-DC] points (D$-$C). The corresponding downstream changes are [DEX-BA], [DEX-CB], and [DEX-DC] EX points. Thus, [RESULT-C4: state whether the observed gain is primarily attributable to KB-conditioned generation, the additional refinement stage, knowledge grounding within the loop, or a combination]. The C$-$B comparison also tests whether ungrounded self-correction is sufficient, a distinction found important in prior tool-interactive refinement work~\cite{gou2024critic}.

The class-level results further show that [RESULT-C5: identify the error classes reduced by generic refinement and those additionally reduced by knowledge grounding]. In particular, [REPORT II-1, II-2, AND III-1 CHANGES FROM C TO D], linking the aggregate gain to the error classes whose evidence resides in the knowledge base. The complete per-class breakdown and secondary rationale and knowledge-group ablations are reported in Appendix~\ref{app:verification-ablation-details}.



\subsubsection{External Controls}

\paragraph{Teacher Control.}
\label{subsec:teacher-control}

The preceding ablation separates the contributions of knowledge grounding and refinement within the pipeline. We now test whether the downstream benefit depends on a much larger teacher. We compare direct prompting, simple chain-of-thought (CoT) prompting, and SemanticAgent under the same Qwen3-235B-A22B-Instruct driver, and then compare direct prompting with SemanticAgent when the driver has the same size and architecture as the downstream student ([STUDENT-EQ]). All arms use matched schemas, task distributions, candidate budgets, final sample counts, student initialization, and fine-tuning settings. The primary readout is the verifiable semantic error rate of the resulting corpus; downstream execution accuracy on [BENCH-E5] is the secondary readout.

\begin{table*}[!t]
\centering
\caption{Teacher control under matched synthesis and training conditions. The upper block holds the driver fixed at Qwen3-235B-A22B; the lower block uses a driver with the same size and architecture as the downstream student ([STUDENT-EQ]). The same-size comparison removes the teacher--student capacity gap but does not rule out self-training or self-distillation. All values are placeholders.}
\label{tab:teacher-control}
\small
\setlength{\tabcolsep}{4pt}
\renewcommand{\arraystretch}{1.15}
\begin{tabular}{llccccc}
\toprule
\textbf{Driver} & \textbf{Protocol} & \textbf{KB construction} & \textbf{Refinement} & \textbf{Final samples} & \textbf{Error rate} & \textbf{EX on [BENCH-E5]} \\
\midrule
Qwen3-235B & Direct & no & no & [N] & [ERR-D235] & [EX-D235] \\
Qwen3-235B & Simple CoT & no & no & [N] & [ERR-C235] & [EX-C235] \\
Qwen3-235B & SemanticAgent & yes & yes & [N] & [ERR-S235] & [EX-S235] \\
\midrule
\texttt{STUDENT-EQ} & Direct & no & no & [N] & [ERR-DSS] & [EX-DSS] \\
\texttt{STUDENT-EQ} & SemanticAgent & independently built & yes & [N] & [ERR-SSS] & [EX-SSS] \\
\bottomrule
\end{tabular}
\end{table*}

With Qwen3-235B held fixed, SemanticAgent yields an error rate of [ERR-S235], compared with [ERR-C235] for Simple CoT and [ERR-D235] for direct prompting; the corresponding downstream scores on [BENCH-E5] are [EX-S235], [EX-C235], and [EX-D235], respectively (Table~\ref{tab:teacher-control}). Thus, [RESULT-Q1: state whether SemanticAgent improves over both simpler protocols, improves corpus quality without improving EX, or is comparable to Simple CoT]. When the driver and downstream student share the same size and architecture, the error rate changes from [ERR-DSS] to [ERR-SSS] and EX changes from [EX-DSS] to [EX-SSS]. This comparison removes the teacher--student capacity gap but not self-training effects; taken together, the two controls test whether SemanticAgent retains an advantage over simpler synthesis protocols under a fixed strong teacher and under a same-size driver, rather than establishing that its gains are independent of all knowledge transfer.



\paragraph{Generator Control under a Fixed Question Set.}
\label{subsec:generator-control}

The teacher controls test the effect of the synthesis driver. We additionally compare MAC-SQL~\cite{wang2025macsql} and MAGIC~\cite{askari2025magic} on a fixed set of questions synthesized by SemanticAgent. Each framework receives the same questions and database schemas, runs its original generation and correction procedure without the SemanticAgent knowledge base, and retains its own intermediate record. The resulting corpora are evaluated with the same contamination, sampling, and downstream fine-tuning protocol. EMLC~\cite{lei2026emlc} is discussed in Related Work but is not included because no public implementation is available.

\begin{table}[!t]
\centering
\caption{Question-controlled comparison of correction-capable SQL generators. Values marked with $[\cdot]$ are placeholders.}
\label{tab:generator-control}
\small
\setlength{\tabcolsep}{4pt}
\begin{tabular}{lcccc}
\toprule
\textbf{Generator} & \textbf{Sampled} & \textbf{SER (\%)} & \textbf{Error rate (\%)} & \textbf{EX} \\
\midrule
MAC-SQL & $[\cdot]$ & $[\cdot]$ & $[\cdot]$ & $[\cdot]$ \\
MAGIC & $[\cdot]$ & $[\cdot]$ & $[\cdot]$ & $[\cdot]$ \\
SemanticAgent & $[\cdot]$ & $[\cdot]$ & $[\cdot]$ & $[\cdot]$ \\
\bottomrule
\end{tabular}
\end{table}

The comparison isolates the SQL-generation and correction stage under a fixed question set; its semantic-quality and downstream results are interpreted together after the measurements are available.

\subsection{Scaling Behavior (RQ4)}

We examine how downstream performance changes with the synthetic data budget under the same backbone and training protocol. The analysis asks whether SemanticAgent's advantage appears in the low-resource regime and whether it remains as the budget increases.

\begin{figure}[!t]
    \centering
    \includegraphics[width=\columnwidth]{performance_scaling_law.pdf}
    \caption{Downstream performance as a function of the synthetic data budget.}
    \label{fig:scaling_law}
\end{figure}

Figure~\ref{fig:scaling_law} shows that SemanticAgent [SCALING-RESULT]. Its advantage [LOW-RESOURCE-RESULT] and [LARGE-BUDGET-RESULT]. [SCALING-INTERPRETATION].

\subsection{Controlled Reanalysis of Semantic Alignment on BIRD}
\label{subsec:bird-sa-reanalysis}

On BIRD, OmniSQL has both a higher SER and a higher SA than SemanticAgent (94.83\% versus 92.14\% SER; 97.41\% versus 91.31\% SA). This pattern cannot be explained solely by the claim that a lower SER leaves only simpler queries in the executable subset. We therefore control the difficulty composition of the executable subsets and compare SA with human judgments on the same evaluated pairs.

\paragraph{Difficulty-controlled comparison.}
We divide executable BIRD queries into the four AST-based complexity levels used in the sampling protocol and sample the same number from every method within each level. The matched-sample size is fixed by the smallest available stratum, and the construction rule is fixed before inspecting the results. We also draw [N-RAND] random subsets of the same size; their maximum SA variation of [RAND-THRESH] points provides a reference for small differences.

\begin{table*}[!t]
\centering
\caption{Semantic Alignment on BIRD under difficulty control. Raw SA is computed over each method's executable subset; Matched SA is computed over subsets with identical difficulty composition. Values marked with $[\cdot]$ are placeholders.}
\label{tab:bird-sa-matched}
\small
\setlength{\tabcolsep}{4pt}
\resizebox{\textwidth}{!}{%
\begin{tabular}{lccccccc}
\toprule
\textbf{Method} & \textbf{SER (\%)} & \textbf{Raw SA (\%)} & \textbf{Easy} & \textbf{Med.} & \textbf{Hard} & \textbf{XHard} & \textbf{Matched SA (\%)} \\
\midrule
CodeS & 91.37 & 95.12 & $[\cdot]$ & $[\cdot]$ & $[\cdot]$ & $[\cdot]$ & $[\cdot]$ \\
SynQL & 90.12 & 92.16 & $[\cdot]$ & $[\cdot]$ & $[\cdot]$ & $[\cdot]$ & $[\cdot]$ \\
OmniSQL & 94.83 & 97.41 & $[\cdot]$ & $[\cdot]$ & $[\cdot]$ & $[\cdot]$ & $[\cdot]$ \\
SemanticAgent & 92.14 & 91.31 & $[\cdot]$ & $[\cdot]$ & $[\cdot]$ & $[\cdot]$ & $[\cdot]$ \\
\bottomrule
\end{tabular}%
}
\end{table*}

Under matched difficulty, the SA gap changes from [GAP-RAW] to [GAP-MATCHED] points, while [WITHIN-LEVEL-DIFFERENCES]. The result [supports/does not support] the explanation that the raw BIRD gap is caused mainly by difficulty composition.

\paragraph{Comparison with human judgments.}
On the shared executable BIRD samples, pairs with and without a human-confirmed violation differ by [SA-DIFF] SA points. Among pairs in the highest SA quartile, [SA-FP]\% still receive a human error label. The SA ranking [agrees with/reverses] the ranking based on the human-evaluated error rates in Section~\ref{subsec:direct-semantic-evaluation}. SA is therefore read as an auxiliary indicator, while direct human evaluation remains the primary evidence for corpus semantic quality.

Benchmark gold-reference noise affects downstream EX and TS but not this comparison, because SA uses each method's own question--SQL pairs and does not consult benchmark gold queries.

\section{Conclusion}

In this paper, we identify a key limitation of existing text-to-SQL synthesis pipelines: execution success does not guarantee semantic correctness. Synthetic samples may execute successfully while still violating schema semantics or domain-specific constraints, which can reduce synthetic data quality and limit downstream performance. To address this issue, we propose \textbf{SemanticAgent}, a knowledge-guided synthesis framework built on a structured knowledge base derived from database schemas and data. This knowledge base supports both question--SQL generation and sample verification. By validating generated pairs against schema semantics and domain-specific constraints, SemanticAgent reduces semantically inconsistent yet executable samples and improves the quality of synthesized training data.

Under matched backbones and training protocols, models fine-tuned on SemanticAgent-generated data consistently outperform those trained on prior synthetic corpora. The gains are largest on BIRD and Spider~2.0-SQLite, where stronger schema grounding and domain-specific reasoning are required. These results suggest that knowledge-guided verification is especially beneficial in semantically demanding text-to-SQL settings.

Several limitations remain. First, offline data synthesis is computationally expensive. Constructing the full training corpus requires approximately 2{,}880 GPU-hours, with verification and refinement accounting for 74.3\% of the total cost. Second, the current pipeline relies on a single teacher model, so bias in that model or contamination from its pretraining may propagate to the synthesized data. Third, our evaluation is limited to public academic benchmarks and may not fully capture the schema noise, SQL dialect variation, and deployment constraints of proprietary enterprise environments. Finally, we primarily use execution accuracy as the evaluation metric, which may still reward semantically incorrect SQL when it happens to produce the correct output.

Future work should focus on making verification and refinement more efficient, reducing dependence on a single teacher model, and evaluating knowledge-guided synthesis under noisier schemas and more realistic deployment settings. More broadly, our findings suggest that executability alone is insufficient for reliable synthetic supervision in text-to-SQL, and that explicit semantic verification is a promising direction for improving downstream performance after fine-tuning.


\section*{Acknowledgments}

\section*{Declaration of competing interest}
The authors declare that they have no known competing financial interests or personal relationships that could have appeared to influence the work reported in this paper.

\section*{Declaration of generative AI and AI-assisted technologies in the writing process}
During the preparation of this work the authors used ChatGPT in order to improve writing. After using this tool/service, the authors reviewed and edited the content as needed and take full responsibility for the content of the publication.

\bibliographystyle{IEEEtran}
\bibliography{ieee-refs}





\clearpage
\appendices
\onecolumn
\subsection{ER Analysis}
\label{app:er-analysis}

% To fit Elsevier 1p single-column layout, the ER content is split into three compact tables with fixed column widths (p{...}), tight spacing, and small font for readability.

\begin{table*}[!ht]
\centering
\small
\setlength{\tabcolsep}{4pt}
\renewcommand{\arraystretch}{1.08}
\caption{Entity identifiers, categorical descriptors, and quantitative indicators for text-to-SQL synthesis.}
\label{tab:er-indicators}
\begin{tabular}{@{}llp{0.62\linewidth}@{}}
\toprule
\textbf{Entity} & \textbf{Type} & \textbf{Description} \\
\midrule
School & CD & School category; district; county; charter status \\
School & CD & Geographic coordinates (latitude, longitude) \\
School & CD & Magnet program; operational status (active) \\
School & QI & Quantitative indicators tracked under AP/SE entities \\
\midrule
SE & CD & IRC category code (e.g., \texttt{1}) \\
SE & CD & NSLP provision status \\
SE & CD & CALPADS certification status \\
SE & QI & FRPM eligibility rate (ages 5--17) \\
\midrule
AP & CD & Subject areas (reading, mathematics, writing) \\
AP & QI & Mean reading score; mean mathematics score \\
AP & QI & Mean writing score; count of scores (\(\geq 1500\)) \\
AP & QI & Participation rate (\texttt{NumTestTakers}/\texttt{Grade12Enrollment}) \\
\midrule
ADM & CD & Opening date; closing date; charter ID number \\
ADM & CD & Virtual campus indicator; funding type \\
ADM & QI & District of Choice (DOC), e.g., \texttt{code 00} \\
\bottomrule
\end{tabular}
\end{table*}

\begin{table*}[!ht]
\centering
\footnotesize
\setlength{\tabcolsep}{3pt}

\begin{minipage}[t]{0.30\textwidth}
    \centering
    \caption{Entity functional roles.}
    \label{tab:er-core}
    \begin{tabular}{@{}ll@{}}
        \toprule
        \textbf{Entity} & \textbf{Role} \\
        \midrule
        SCH & Primary entity \\
        SE  & Domain attribute \\
        AP  & Domain attribute \\
        ADM & Metadata entity \\
        \bottomrule
    \end{tabular}
\end{minipage}
\hfill
\begin{minipage}[t]{0.35\textwidth}
    \centering
    \caption{Source tables by entity.}
    \label{tab:er-sources}
    \begin{tabular}{@{}lp{0.72\linewidth}@{}}
        \toprule
        \textbf{Entity} & \textbf{Source Tables} \\
        \midrule
        SCH & \texttt{schools}; \texttt{satscores}; \texttt{frpm} \\
        SE  & \texttt{frpm}; \texttt{cdscode}; \texttt{AY} \\
        AP  & \texttt{satscores}; \texttt{cdscode} \\
        ADM & \texttt{schools}; \texttt{cdscode} \\
        \bottomrule
    \end{tabular}
\end{minipage}

\begin{minipage}{0.9\textwidth}
\footnotesize
\vspace{1mm}
\textit{Note:} 
\noindent Entity abbreviations used throughout this analysis: School (\textbf{SCH}), Socioeconomic eligibility (\textbf{SE}), Academic performance (\textbf{AP}), \texttt{Academic year} (\texttt{AY}), Administrative metadata (\textbf{ADM}), Categorical descriptors (\textbf{CD}) and Quantitative indicators (\textbf{QI}).
\end{minipage}
\end{table*}


\clearpage
\subsection{Synthetic data format with structured reasoning}
\label{app:training_format}

\lstdefinestyle{jsonStyle}{
    basicstyle=\scriptsize\ttfamily,
    numbers=none,
    breaklines=true,
    breakatwhitespace=true,
    tabsize=2,
    showstringspaces=false,
    columns=flexible,
    keepspaces=true,
    frame=none
}

\begin{figure*}[!ht]
\centering
\begin{tcolorbox}[
  width=.7\linewidth,
  enhanced,
  boxrule=0.5pt,
  colback=AliceBlue,
  colframe=Navy,
  colbacktitle=Navy,
  fonttitle=\bfseries\color{white}, 
  title=Synthetic data format with structured reasoning,
  arc=2mm,
  attach boxed title to top left={yshift=-2.5mm, xshift=3mm},
  boxed title style={arc=1mm},
  boxsep=5pt,
  left=8pt,
  right=8pt,
  top=6mm,
]
\begin{lstlisting}[style=jsonStyle]
{
  "question": "Which school has the highest average SAT math score?",
  "think": {
    "focus": "Identify single school with MAX(AvgScrMath)",
    "metadata": {
      "main_scenario": "ranking_query",
      "sub_scenario": "max_min_query",
      "complexity_level": 1,
      "use_case": "ranking_analysis"
    },
    "table_selection": {
      "tables_used": ["satscores"],
      "reasoning": "Contains SAT columns for extreme value. Single-table compliant with Level 1"
    },
    "column_selection": {
      "columns_used": [
        {
          "name": "sname",
          "type": "TEXT",
          "operation": "SELECT",
          "purpose": "Identify school with extreme score"
        },
        {
          "name": "AvgScrMath",
          "type": "INTEGER",
          "operation": "ORDER BY DESC",
          "purpose": "Find maximum via sorting"
        }
      ]
    },
    "sql_strategy": {
      "operations": ["SELECT", "ORDER BY DESC", "LIMIT"],
      "approach": "ORDER BY + LIMIT for extreme value, no aggregation",
      "no_need": ["JOIN", "GROUP BY", "MAX()"]
    },
    "expected_output": "Single row: school name with highest AvgScrMath"
  },
  "answer": "SELECT sname FROM satscores ORDER BY AvgScrMath DESC LIMIT 1;"
}
\end{lstlisting}
\end{tcolorbox}
\caption{Synthetic data format with structured reasoning.}
\label{fig:training_data_format}
\end{figure*}

\clearpage

\subsection{Prompts}
\label{app:prompts}

\begin{figure*}[!ht]
\centering
\begin{tcolorbox}[
  enhanced,
  boxrule=0.5pt,
  colback=AliceBlue,
  colframe=Navy,
  colbacktitle=Navy,
  fonttitle=\bfseries\color{white}, 
  title=Prompt for SQL Complexity Classification,
  arc=2mm,
  attach boxed title to top left={yshift=-2.5mm, xshift=3mm},
  boxed title style={arc=1mm},
  boxsep=5pt,
  left=8pt,
  right=8pt,
  top=6mm,
]
\scriptsize\ttfamily  
You are an SQL complexity evaluator. Classify the given SQL query into one of four levels based on the following criteria. Output only the level number and reasoning.

\vspace{1.5ex}
\textbf{Classification Criteria:}

\vspace{0.5ex}
\textbf{Level 1: Simple}\\
Single table, basic filtering and sorting (WHERE, ORDER BY, LIMIT, DISTINCT)

\vspace{0.5ex}
\textbf{Level 2: Moderate}\\
Multi-table joins (max 3 tables, 2 JOINs) with aggregation (GROUP BY, HAVING, COUNT/SUM/AVG)

\vspace{0.5ex}
\textbf{Level 3: Complex}\\
Advanced logic: subqueries (IN, EXISTS, correlated), CASE WHEN, UNION, or 4+ tables

\vspace{0.5ex}
\textbf{Level 4: Expert}\\
Must use: Window functions (ROW\_NUMBER, RANK, LAG, LEAD), CTEs (WITH), or recursive queries

\vspace{1.5ex}
\textbf{Decision Rule:}\\
Check Level 4 features first $\rightarrow$ Level 3 $\rightarrow$ Level 2 $\rightarrow$ else Level 1

\vspace{1.5ex}
\textbf{Output JSON:}\\
\{``level'': ``number'', ``reasoning'': ``brief explanation''\}

\vspace{1.5ex}
\textbf{SQL Query:}\\
\{sql\_query\}
\end{tcolorbox}
\caption{The prompt used for SQL complexity classification. The model classifies queries into four levels based on their features and constraints.}
\label{fig:prompt_sql_complexity}
\end{figure*}

\begin{figure*}[!ht]
\centering
\begin{tcolorbox}[
  enhanced,
  boxrule=0.5pt,
  colback=AliceBlue,
  colframe=Navy,
  colbacktitle=Navy,
  fonttitle=\bfseries\color{white}, 
  title=Prompt for Question--SQL Semantic Consistency Verification,
  arc=2mm,
  attach boxed title to top left={yshift=-2.5mm, xshift=3mm},
  boxed title style={arc=1mm},
  boxsep=5pt,
  left=8pt,
  right=8pt,
  top=6mm,
]
\scriptsize\ttfamily
You are a semantic consistency evaluator. Determine whether the given SQL query correctly answers the natural language question based on the provided database schema. Output 1 if consistent, 0 if inconsistent.\\[2ex]

\textbf{Evaluation Criteria:}\\[1ex]
\textbf{Consistent (1):} The SQL query accurately retrieves the information requested in the question, using correct tables, columns, joins, filters, and aggregations according to the schema.\\[1ex]

\textbf{Inconsistent (0):} The SQL query has one or more issues:\\
- Targets wrong tables or columns\\
- Uses incorrect join conditions or missing necessary joins\\
- Applies wrong filters, aggregations, or ordering\\
- Returns irrelevant or incomplete data for the question\\[2ex]

\textbf{Verification Steps:}\\
1. Parse the question to identify required information\\
2. Check if SQL uses correct tables/columns from schema\\
3. Verify joins, filters, and aggregations match the question intent\\
4. Confirm the output answers the question completely\\[2ex]

\textbf{Output:} \{``label'': 0 or 1, ``reasoning'': brief explanation\}\\[2ex]

\textbf{Database Schema:}\\
\{schema\}\\[2ex]

\textbf{Natural Language Question:}\\
\{question\}\\[2ex]

\textbf{SQL Query:}\\
\{sql\_query\}
\end{tcolorbox}
\caption{The prompt used for Question--SQL semantic consistency verification. The model evaluates whether the SQL query correctly answers the question given the database schema.}
\label{fig:prompt_consistency}
\end{figure*}

\subsection{Semantic Error Taxonomy and Prompts}
\label{app:taxonomy}

This appendix provides the full description of each class in Table~\ref{tab:semantic-taxonomy}, together with the prompt templates used for semantic checking and correction.

\subsubsection{Class Descriptions}

Table~\ref{tab:taxonomy-details} restates each class with a checkable description. Placeholders denote schema components: [A] column, [C] condition, [J] join, [V] value, [R] domain rule.

\begin{table*}[!t]
\centering
\caption{Full description of the semantic error classes.}
\label{tab:taxonomy-details}
\small
\setlength{\tabcolsep}{6pt}
\renewcommand{\arraystretch}{1.15}
\begin{tabular}{clll}
\toprule
\textbf{ID} & \textbf{Class} & \textbf{Description} & \textbf{Evidence} \\
\midrule
\multicolumn{4}{l}{\emph{Group I: intent mismatches}} \\
I-1 & Task mismatch & The task $s$ performs differs from the task $q$ requests. & $q$, $s$ \\
I-2 & Condition mismatch & The condition [C] expressed in $q$ is missing or altered in $s$. & $q$, $s$ \\
I-3 & Aggregation/granularity mismatch & The aggregation or granularity differs from the question. & $q$, $s$ \\
I-4 & Ordering/result-form mismatch & The ordering or result form differs from the question. & $q$, $s$ \\
\midrule
\multicolumn{4}{l}{\emph{Group II: element misuse}} \\
II-1 & Field-role misuse & The operation applied to column [A] contradicts its real-world meaning. & $\mathcal{K}_3$, $\mathcal{K}_4$ \\
II-2 & Join misuse & The join [J] expresses a relationship that does not exist. & $\mathcal{K}_6$ \\
\midrule
\multicolumn{4}{l}{\emph{Group III: domain-rule violations}} \\
III-1 & Domain-rule violation & The value [V] or filter violates a domain rule [R]. & $\mathcal{K}_2$, $\mathcal{K}_5$, $\mathcal{K}_6$ \\
\bottomrule
\end{tabular}
\end{table*}

\subsubsection{Diagnostic Prompt}
\label{app:taxonomy:diagnose}

The semantic checks described in Section~\ref{subsec:refinement} are performed with the following prompt template; the knowledge entries retrieved for the involved schema elements are inserted at the designated position.

% TODO: 以下为骨架示意，请以实现中的真实 prompt 替换（保留输出格式约束）。
\begin{quote}
\small
\texttt{[System]} You are a semantic verifier for text-to-SQL data synthesis. \\
\texttt{[Input]} Question $q$; SQL query $s$; rationale $r$; knowledge entries from $\mathcal{K}$: \texttt{<retrieved entries>}. \\
\texttt{[Task]} Examine the pair in order: (1) Is the question valid in the target domain? (2) Does $s$ capture the intent of $q$---its task, conditions, aggregation, granularity, and result form? (3) Are the involved columns, values, and joins used consistently with their real-world meanings? \\
\texttt{[Output]} For each violation, report: class (one of the seven classes of Table~\ref{tab:taxonomy-details}), affected artifact ($q$/$s$/$r$), location, and the violated condition.
\end{quote}

\subsubsection{Correction Prompt}
\label{app:taxonomy:correct}

The \textsc{Correct} function edits the SQL query and the rationale with the following prompt template; the question is never modified.

% TODO: 以下为骨架示意，请以实现中的真实 prompt 替换（保留"不得修改问题"的约束）。
\begin{quote}
\small
\texttt{[System]} You are a semantic corrector for text-to-SQL data synthesis. \\
\texttt{[Input]} Question $q$ (fixed, never modify); SQL query $s$; rationale $r$; diagnosed violations with their evidence \texttt{<diagnosis--evidence pairs>}; execution and schema errors \texttt{<B entries>}. \\
\texttt{[Task]} Apply local edits to $s$ and $r$ to resolve each violation, conditioned on its evidence. Do not regenerate the query or modify the question. \\
\texttt{[Output]} The corrected query and rationale, together with the edits applied.
\end{quote}

\subsection{Per-Class Diagnosis Results}
\label{app:diagnosis-by-type}

Table~\ref{tab:recall-by-type} in the main text reports class-level recall for the three knowledge-grounded classes and the merged intent-level classes. This appendix reports all seven classes individually, with the number of drafts having at least one human-confirmed violation of each class.

\begin{table}[!t]
\centering
\caption{Full seven-class breakdown of the class-level recall of the two configurations. Denominators are as in Table~\ref{tab:recall-by-type}. Values marked with $[\cdot]$ are placeholders to be replaced by the measured results.}
\label{tab:recall-by-type-full}
\small
\setlength{\tabcolsep}{4pt}
\begin{tabular}{lrrr}
\toprule
\textbf{Class} & \textbf{$n$} & \textbf{Generic critic} & \textbf{KB-cond.} \\
\midrule
I-1 Task mismatch & [N] & [R] & [R] \\
I-2 Condition mismatch & [N] & [R] & [R] \\
I-3 Aggregation/granularity mismatch & [N] & [R] & [R] \\
I-4 Ordering/result-form mismatch & [N] & [R] & [R] \\
II-1 Field-role misuse & [N] & [R] & [R] \\
II-2 Join misuse & [N] & [R] & [R] \\
III-1 Domain-rule violation & [N] & [R] & [R] \\
\bottomrule
\end{tabular}
\end{table}

\subsection{Representative Refinement Cases}
\label{app:refinement-cases}

This appendix illustrates four refinement outcomes reported in Section~\ref{subsec:paired-refinement}. The examples are selected after annotation and are intended to show the complete transition from the initial draft to the final disposition, rather than to represent an average case. All identifiers and values remain placeholders until the corresponding logs and adjudicated annotations are finalized.

\subsubsection{Successful repair}
\label{app:refinement-cases:repair}

\begin{description}
\item[Case ID:] [CASE-REPAIR]
\item[Initial triple:] Question [Q-REPAIR]; SQL [SQL-REPAIR]; rationale [R-REPAIR].
\item[Initial labels:] [INITIAL-LABELS-REPAIR].
\item[Refinement trace:] Diagnosis [DIAG-REPAIR-1] with evidence [EVIDENCE-REPAIR-1]; correction [EDIT-REPAIR-1]; re-verification [RECHECK-REPAIR].
\item[Final assessment:] [FINAL-LABELS-REPAIR]; disposition: repaired.
\end{description}

\subsubsection{Mis-repair}
\label{app:refinement-cases:misrepair}

\begin{description}
\item[Case ID:] [CASE-MISREPAIR]
\item[Initial triple:] Question [Q-MISREPAIR]; SQL [SQL-MISREPAIR]; rationale [R-MISREPAIR].
\item[Initial labels:] [INITIAL-LABELS-MISREPAIR].
\item[Refinement trace:] Diagnosis [DIAG-MISREPAIR] and edit [EDIT-MISREPAIR].
\item[Final assessment:] [FINAL-LABELS-MISREPAIR], including the newly introduced violation [NEW-LABEL-MISREPAIR]; disposition: mis-repaired.
\end{description}

\subsubsection{Rejection at the iteration limit}
\label{app:refinement-cases:rejection}

\begin{description}
\item[Case ID:] [CASE-REJECT]
\item[Initial triple:] Question [Q-REJECT]; SQL [SQL-REJECT]; rationale [R-REJECT].
\item[Refinement trace:] Iteration diagnoses and edits [TRACE-REJECT].
\item[Final disposition:] The violation [RESIDUAL-LABEL-REJECT] remains at the limit [MAX-ITER]; the draft is rejected and does not enter the final corpus.
\end{description}

\subsubsection{Damage to an initially clean draft}
\label{app:refinement-cases:damage}

\begin{description}
\item[Case ID:] [CASE-DAMAGE]
\item[Initial triple:] Question [Q-DAMAGE]; SQL [SQL-DAMAGE]; rationale [R-DAMAGE].
\item[Initial labels:] no human-confirmed semantic violation.
\item[Refinement trace:] Diagnosis [DIAG-DAMAGE] followed by edit [EDIT-DAMAGE].
\item[Final assessment:] [FINAL-LABELS-DAMAGE] or non-executable SQL [EXEC-DAMAGE]; disposition: damaged.
\end{description}

\subsection{Supplementary Verification Ablations}
\label{app:verification-ablation-details}

We provide two supplementary analyses of the full SemanticAgent configuration. First, we remove the rationale from either the synthesis-and-verification process or the downstream training format, because these operations use rationale in different roles. Second, we remove one knowledge group at a time from both generation and verification. These analyses are secondary to the four-configuration comparison in Section~\ref{subsec:ablation}.

\begin{table}[!t]
\centering
\caption{Supplementary rationale and loop ablations. Values marked with $[\cdot]$ are placeholders.}
\label{tab:ablation-rationale-loop}
\small
\setlength{\tabcolsep}{4pt}
\renewcommand{\arraystretch}{1.15}
\begin{tabular}{lccc}
\toprule
\textbf{Configuration} & \textbf{Error rate} & \textbf{Repair success} & \textbf{Damage} \\
\midrule
Full (reference) & [ERR-FULL] & [REP-FULL] & [DMG-FULL] \\
w/o rationale in synthesis/verification & [ERR-NOR-P] & [REP-NOR-P] & [DMG-NOR-P] \\
w/o rationale in training format & [ERR-NOR-T] & --- & --- \\
KB-conditioned generation + generic refinement & [ERR-CGENB] & [REP-CGENB] & [DMG-CGENB] \\
\bottomrule
\end{tabular}
\end{table}

Removing [RATIONALE ROLE] changes the error rate by [DELTA-RATIONALE] points, whereas combining KB-conditioned generation with generic refinement changes it by [DELTA-CGENB] points relative to the full pipeline. We interpret these differences as secondary configuration effects rather than as additional estimates of the main KB-versus-refinement contrast.

\begin{table*}[!t]
\centering
\caption{Knowledge-group removal. One group is removed from both generation and validation, and the corpus is regenerated under the same budget. Values marked with $[\cdot]$ are placeholders.}
\label{tab:knowledge-split}
\small
\setlength{\tabcolsep}{4pt}
\renewcommand{\arraystretch}{1.15}
\begin{tabular}{lccc}
\toprule
\textbf{Removed group} & \textbf{Error rate} & \textbf{Expected class} & \textbf{Observed change} \\
\midrule
None (full reference) & [ERR-D] & --- & --- \\
$\mathcal{K}_2,\mathcal{K}_5$ & [ERR-NODT] & III-1 (partial) & [OBS-NODT] \\
$\mathcal{K}_3,\mathcal{K}_4$ & [ERR-NOFC] & II-1 & [OBS-NOFC] \\
$\mathcal{K}_6$ & [ERR-NOREL] & II-2; III-1 (partial) & [OBS-NOREL] \\
\bottomrule
\end{tabular}
\end{table*}

Removing $\mathcal{K}_3,\mathcal{K}_4$ restores [OBS-NOFC], while removing $\mathcal{K}_6$ restores [OBS-NOREL]. These changes [support/do not support] the expected correspondence between knowledge groups and semantic error classes. This analysis measures dependence on the knowledge groups; the correctness of the entries themselves is evaluated separately in Section~\ref{subsec:kb-audit}.

\clearpage


\subsection{Annotation Guideline}
\label{app:annotation-guideline}

This appendix provides the full annotation guideline used in the human evaluation of Section~\ref{subsec:human-evaluation}: the judgment sequence, the decision points for each condition of Definition~1, and worked examples for all seven classes of Table~\ref{tab:semantic-taxonomy}. [CONTENT-TBD: paste the finalized guideline.]

\subsection{Per-Database Audit Results}
\label{app:kb-stages}

Table~\ref{tab:kb-audit-perdb} reports the entry-level audit of Section~\ref{subsec:kb-audit} broken down by source database. Values marked with $[\cdot]$ are placeholders until the audit logs are finalized.

\begin{table}[!t]
\centering
\caption{Per-database distribution of audit labels for the LLM-inferred knowledge layers. Values marked with $[\cdot]$ are placeholders.}
\label{tab:kb-audit-perdb}
\small
\setlength{\tabcolsep}{3pt}
\resizebox{\columnwidth}{!}{%
\begin{tabular}{lccccc}
\toprule
\textbf{Database} & \textbf{Sampled} & \textbf{Correct and supported} & \textbf{Incorrect} & \textbf{Plausible but unsupported} & \textbf{Not assessable} \\
\midrule
DB-1 & [N] & [n (\%)] & [n (\%)] & [n (\%)] & [n (\%)] \\
DB-2 & [N] & [n (\%)] & [n (\%)] & [n (\%)] & [n (\%)] \\
\bottomrule
\end{tabular}%
}
\end{table}

% \section{Biography Section}
% If you have an EPS/PDF photo (graphicx package needed), extra braces are
%  needed around the contents of the optional argument to biography to prevent
%  the LaTeX parser from getting confused when it sees the complicated
%  $\backslash${\tt{includegraphics}} command within an optional argument. (You can create
%  your own custom macro containing the $\backslash${\tt{includegraphics}} command to make things
%  simpler here.)
 
% \vspace{11pt}

% \bf{If you include a photo:}\vspace{-33pt}
% \begin{IEEEbiography}[{\includegraphics[width=1in,height=1.25in,clip,keepaspectratio]{fig1}}]{Michael Shell}
% Use $\backslash${\tt{begin\{IEEEbiography\}}} and then for the 1st argument use $\backslash${\tt{includegraphics}} to declare and link the author photo.
% Use the author name as the 3rd argument followed by the biography text.
% \end{IEEEbiography}

% \vspace{11pt}

% \bf{If you will not include a photo:}\vspace{-33pt}
% \begin{IEEEbiographynophoto}{John Doe}
% Use $\backslash${\tt{begin\{IEEEbiographynophoto\}}} and the author name as the argument followed by the biography text.
% \end{IEEEbiographynophoto}




\vfill

\end{document}


